%% 第26回編集委員会（物理第4巻）：組 電W8
%% p4-15.tex  電-15　光も電波も同じ波：電磁波（★2〜5・3段）
%% 中身は計画書 tools/phys4-plan.md の B「電-15」。公式の持ち主は A（A58・A59・A85〜A88、A19 の電磁波の利用）。
%% 数値はすべて Python で計算した：checks/物理第4巻/電-15/calc_p4_15.py（出力 calc_p4_15_out.txt）。
%% 記号（この章。巻全体は計画書 C-1・第26回 Y-4）：
%%   光速 c（第4部。W-4）、周波数 f（波動の本の振動数と同じ量：Y-8）、角周波数 ω、波長 λ、波数 k（クーロンの定数は 1/(4πε0) と書き、k と書かない）、
%%   電場 E（振幅 E_0）、磁場 B（振幅 B_0）、磁場の強さ H（振幅 H_0）、波のインピーダンス η（真空 η_0）、ポインティングベクトル S、
%%   強さ I〔W/m²〕（電流は記号にしない：計画書 C-1。波-7 と同じ）、面積 A（第4部：Y-4）、放射圧 p（運動量は記号にしない）、
%%   電荷 q、加速度 a、電子のゆれの振幅 x_0（波-13 の A。この章では A は面積なので x_0 と書く）、
%%   ゆれの向きと見る向きの角 θ、距離 r、出す電磁波の仕事率 P（分極 P は使わない）、屈折率 n、比誘電率 ε_r、比透磁率 μ_r、
%%   入射角 θ_1・屈折角 θ_2、ブルースター角 θ_B（カードは波-13）、電場の反射係数 r_E（波-7 の r_押 と同じ式：X-3）、
%%   反射率 R（電気抵抗は出ない。放射抵抗は R_r と書く）、導電率 σ、表皮深さ δ（持ち主は電-8 ⑫）、
%%   網の穴の直径 d、網の厚さ L（インダクタンスは出ない）、遮断波長 λ_c、時刻 t。
%% 公式の持ち主：c＝fλ・電磁波の分類・横波・電磁波の発生（A58・A59：高校物理）、波動方程式と c＝1/√(ε0μ0)・E＝cB・η0（A85）、
%%   放射圧（A86）、ラーモアの式（A87）、屈折率 n＝√(εrμr)・反射率・ブルースター角の理由（A88）はこの章。
%%   マクスウェル方程式・変位電流・ポインティングベクトルは電-14、表皮深さは電-8 ⑫、境界条件は電-3 ⑨・電-7 ⑦（静的な場合）、
%%   特性インピーダンスと電圧の反射係数は電-13 ⑦⑧、逆2乗の法則は波-7 ②、強制振動は波-2 ④、屈折の法則は波-12 ③。
%%   計算が済んでいる所（8-14⑤ の波動方程式、8-12⑧（S15）の rot rot、8-13⑩（S11）の波数）は数学の章を指す（Y-9）。
%% 前の巻との約束（計画書 A-6）：P12・P42（空はなぜ青い）は⑧、P19（熱-13：光が電気と磁気の波）は③④⑥、P25（波-5：媒質なしで進む・波動方程式）は⑥、
%%   P29（波-6：電子レンジ 2.45 GHz・12.2 cm）ははてな？と例題1、P31（波-7：Z∝1/n）は⑦⑨、P35（W-8：出す光の強さは加速度の2乗に比例）は⑧で導いた、
%%   P36（波-13：屈折率の大きい物質での反射で山と谷が入れかわる理由）は⑨、P37（光の電場は進む向きに垂直）は④、
%%   P38（ゆれる電子はゆれの向きに光を出さない・反射率の図 n＝1.5・4 %・56.3°）は⑧⑨と例題8・9、P39 は全体、
%%   P40（波-14：アンテナから電波が出るしくみ、2.4 GHz の λ/4＝3.1 cm）は⑪と例題13、P43（つぎの巻へ）は⑥⑧。
%%   第3巻の記号・数にそろえた：波-13 のグラフ（実線＝紙面に垂直なゆれ、破線＝紙面の中のゆれ）、波-13 の 450・550・700 nm と 5.9 倍・2.6 倍、
%%   波-12 の屈折率の表（水 1.333、ガラス 約1.5、ダイヤモンド 2.42、空気 1.0003）、波-14 の FM 80 MHz・3.75 m、AM 1000 kHz・300 m、Wi-Fi 2.4 GHz・12.5 cm。
%% 出典（第26回 Y-6 2章の16。11月5日に確かめた）：
%%   ・電波の定義（三百万メガヘルツ以下の周波数の電磁波）：電波法 第2条（e-Gov 法令 API 325AC0000000131 で本文を読んだ）。
%%   ・電子レンジ 2,450 MHz±50 MHz（高周波出力の定格 2 kW 以下）：電波法施行規則 第46条の7（e-Gov 法令 API 325M50080000014 で本文を読んだ）。
%%   ・IC カード（誘導式読み書き通信設備）13.56 MHz：同 第44条（同じ API で本文を読んだ）。IH 20.05〜100 kHz：同 第46条の7（計画書 D14）。
%%   ・中波（AM）放送 526.5〜1606.5 kHz：総務省 関東総合通信局「ラジオ放送の概要」（計画書 D16）。
%%   ・FM 放送の周波数 76.1〜94.9 MHz：総務省 信越総合通信局 報道資料（2020年7月17日、ワイドFM の説明。soumu.go.jp で本文を読んだ）。本文は「76〜95 MHz」。
%%   ・地上デジタルテレビ 13〜52 ch（470〜710 MHz）：総務省 報道資料（平成23年8月26日「地上デジタルテレビジョン放送の一部周波数の変更」。本文を読んだ）。
%%   ・電波防護指針（一般環境、6分間平均）1.5〜300 GHz：電界強度 61.4 V/m、磁界強度 0.163 A/m、電力束密度 1 mW/cm²：
%%     総務省 九州総合通信局「電波防護のための基準」（soumu.go.jp/soutsu/kyushu/re/defense.html。本文の表1を読んだ。影響が出る値に約50倍の安全率）、
%%     電波法施行規則 第21条の4・別表第2号の3の3（同じ値。e-Gov で読んだ）。
%%   ・太陽定数 1361 W/m²：IAU 2015 Resolution B3（Prša ほか、arXiv:1605.09788。arxiv.org で読んだ。TSI 1361±0.5 W/m²）。
%%   ・IKAROS の帆：「差し渡し20mの正方形」、厚さ 7.5 μm、2010年に打ち上げ（JAXA プレスリリース 2009年9月9日の資料 PDF を読んだ）。面積 20²/2＝200 m²。
%%   ・高速電波バースト：Lorimer ほか "A bright millisecond radio burst of extragalactic origin"（arXiv:0709.4301、Science 2007）、
%%     銀河系のマグネターから：CHIME/FRB Collaboration（arXiv:2005.10324、Nature 587, 54 (2020)）。どちらも arxiv.org で題と日付を読んだ。
%%   ・ラーモア 1897年 Phil. Mag. 44（DOI 10.1080/14786449708621095 の登録を doi.org で確かめた。題と年は検索で2つの資料が一致）。
%%   ・銅の導電率 5.8×10⁷ S/m：焼きなました銅の国際標準（IACS、1/58 Ω·mm²/m）。検索で copper.org と NIST の資料（1914年の標準）を照らし合わせた。
%%   ・水の比誘電率 80（20 ℃ で 80.2：IAPWS の式。計画書 D17）、水の屈折率 1.333（波-12）。
%%   ・半波長アンテナの放射抵抗 73 Ω：細い半波長アンテナの式 R＝(η0/4π)(γ＋ln2π−Ci(2π)) を Python（scipy）で計算して 73.1 Ω（calc_p4_15.py）。
%%   ・円い穴の遮断波長 1.706 d（ベッセル関数の微分の零点 1.8412 から λc＝πd/1.8412。calc_p4_15.py で計算）。網の穴の直径・厚さは「とします」の仮の値。
%% 漢字：常用漢字だけ（遮へい）。checks/物理第4巻/電-15/kanji_check.py で確かめた。
\tikzset{sxp lab/.style={font=\sffamily\footnotesize, inner sep=1pt},
  sxp E/.style={-{Stealth[length=2.6pt, width=2.4pt]}, line width=0.9pt, bkRed!80!black},
  sxp B/.style={-{Stealth[length=2.6pt, width=2.4pt]}, line width=0.9pt, bkBlue!80!black, dash pattern=on 2.4pt off 1.4pt},
  sxp S/.style={-{Stealth[length=3.2pt, width=3pt]}, line width=1.4pt, bkOrange!85!black},
  sxp ray/.style={-{Stealth[length=3pt, width=2.8pt]}, line width=0.9pt, bkInk},
  sxp dim/.style={{Stealth[length=2.4pt]}-{Stealth[length=2.4pt]}, line width=0.5pt, bkGray!80!black}}

\phchapter{15}{光も電波も同じ波：\\電磁波}{電子レンジの窓の網は、何のためにある？}{\bklevelrange{2}{5}}{光も電波も同じ波：電磁波}
\bkmichishirube{8-8,8-12,8-13,8-14,電-8,電-14,波-7,波-12,波-13,波-14}{電-16,原-2,原-11}

\begin{goal}
\begin{itemize}[label={\color{bkTheme}●}, leftmargin=1.6\zw, itemsep=0pt, topsep=1pt]
\item \bklevel{2}\enspace 電波・光・X 線などが、波長（周波数）の違う同じ電磁波であることを知り、$c=f\lambda$ で波長を計算できる。
\item \bklevel{3}\enspace 電磁波が電場と磁場の直角にゆれる横波で、$E=cB$ の関係をもち、エネルギーと運動量を運ぶことを計算で確かめられる。
\item \bklevel{5}\enspace ゆれる電荷が出す電磁波（ラーモアの式）、物の表面での反射率、金属の網が電波を通さないわけを、マクスウェル方程式から説明できる。
\end{itemize}
\end{goal}

\phkoushiki{%
\phk{電磁波の速さ $c=f\lambda$ と、波長による分類}{高校物理}{ph:電-15:clambda}
\phk{電磁波は横波（電場と磁場が直角にゆれる）}{高校物理}{ph:電-15:yokonami}
\phk{振動する電流（アンテナ）から電磁波が出る}{高校物理}{ph:電-15:hassei}
}

\begin{hatena}
\noindent\begin{minipage}[c]{0.58\linewidth}\raggedright
電子レンジの扉の窓には、細かい穴のあいた金属の網がはりこまれています。網のおかげで、中の食べ物のようすは見えます。ところが、食べ物を温めている電磁波（マイクロ波）は、この網から外へ出てきません。光は通すのに、マイクロ波は通さない。小さな穴のあいた網は、どうやって2つを分けているのでしょう。
\end{minipage}\hfill
\begin{minipage}[c]{0.39\linewidth}\centering
\begin{tikzpicture}[x=1mm, y=1mm]
  \filldraw[fill=bkGray!25, draw=bkInk, line width=0.7pt, rounded corners=1mm] (0,0) rectangle (44,28);
  \filldraw[fill=bkGray!60, draw=bkInk, line width=0.6pt] (3,3) rectangle (31,25);
  \foreach \x in {4.5,6.5,...,29.5} {\foreach \y in {4.5,6.5,...,23.5} {\fill[bkYellow!60] (\x,\y) circle (0.55);}}
  \filldraw[fill=white, draw=bkInk, line width=0.5pt] (34,6) rectangle (41,22);
  \foreach \y in {9,13,17} {\draw[bkInk, line width=0.5pt] (35.5,\y) -- (39.5,\y);}
  \node[sxp lab, anchor=north] at (17,-1) {窓の網（金属）};
\end{tikzpicture}
\end{minipage}
\end{hatena}

\noindent 電-14 では、変わる電場が磁場を作り、変わる磁場が電場を作ることを、4つの式（マクスウェル方程式）にまとめました。この2つがかわるがわる生まれると、電場と磁場のゆれが、空間を次々と伝わっていきます。これが\textbf{電磁波}です。\bkidx{でんじは}{電磁波}ラジオ・テレビ・スマホの電波も、目に見える光も、レントゲンの X 線も、すべて同じ電磁波で、違うのは波長（周波数）だけです。この章では、電磁波の性質を調べ、波動の本（第3巻）で結果だけ使った「光が電磁波であること」「ゆれる電子が出す光の強さ」「反射率の式」を導きます。

%% ---------------- やってみよう・思い出そう ----------------
\begin{yattemiyou}[電波をさえぎる・光のゆれの向きを見る]
\textbf{用意するもの}：スマホ2台（または、スマホと Wi-Fi のルーター）、アルミはく、偏光サングラス（なければ偏光板）
\begin{enumerate}[label={\bkcircled[bkGreen]{\arabic*}}, leftmargin=2.4\zw, itemsep=1pt, topsep=2pt]
\item スマホの画面で、Wi-Fi や電波の強さの表示（アンテナの本数）を見る。
\item スマホをアルミはくで、すきまがないように何重にも包む。もう1台のスマホから電話をかけて、つながるかどうか試す（つながらなくても、包んだまま長く置かない）。
\item 晴れた日に、太陽を背にして、偏光サングラスごしに青空を見る。サングラスを目の前でゆっくり回すと、空の明るさはどう変わるかな？　太陽から $90°$ くらい離れた空で試そう（太陽を直接見ない）。
\end{enumerate}
\textbf{注意}：アルミはくや金属を、電子レンジの中に入れて温めてはいけません。火花が出て危険です。
\end{yattemiyou}

\noindent ②では、アルミはくが電波をはね返し、すきまなく包むと電話がつながりにくくなります（⑩）。③では、空の光が、ある向きにゆれる光を多くふくんでいる（偏光している）ことがわかります（⑧）。

\begin{omoidasou}
\bkrecall{電-14}{マクスウェル方程式：変わる磁場が電場の渦を作り（$\oint\vec E\cdot\dd\vec r=-\dd\Phi_B/\dd t$）、電流と変わる電場が磁場の渦を作る（$\oint\vec B\cdot\dd\vec r=\mu_0I_{\text{中}}+\mu_0\varepsilon_0\,\dd\Phi_E/\dd t$）（電-14 ⑤）。}
\bkrecall{電-14}{エネルギー密度 $\tfrac12\varepsilon_0E^2+B^2/(2\mu_0)$、エネルギーの流れはポインティングベクトル $\vec S=\vec E\times\vec H$（電-14 ⑦）。}
\bkrecall{波-5}{波の速さ＝振動数×波長（$c=f\lambda$）。正弦波 $\sin(\omega t-kx)$ では $\omega=2\pi f$、波数 $k=2\pi/\lambda$（波-5 ②③）。}
\bkrecall{波-13}{散乱：光の電場にゆすられた電子が、同じ振動数の光を出しなおす。出す光の強さが加速度の2乗に比例することは、この章で導く（波-13 ⑧）。}
\bkrecall{波-7}{2つの物の境での反射係数は、インピーダンスの比で決まる。押すもの（力・電圧）で数えた反射係数は $r_{\text{押}}=(Z_2-Z_1)/(Z_1+Z_2)$（波-7 ④）。}
\end{omoidasou}

\noindent この章は3段に分かれています。①〜⑤（目でわかる・手でわかる）は、$c=f\lambda$ と、電-14 の積分形の法則で読めます。⑥〜⑫（式でわかる）では、波動方程式・ゆれる電荷の放射・反射の式・金属の網・アンテナを調べます。面積は、電-14 と同じく $A$ と書きます。

%% ======================== 物理の道具 ========================
\bksection{doguu}{物理の道具}

\phstage{me}{電磁波の家族を並べる}{2}

\bksubsection{1}{電磁波の家族：電波から γ 線まで}

\leavevmode\label{ph:電-15:sec1}電磁波は、波長（周波数）によって、ふるまいも使いみちも大きく変わります。そこで、波長の長いほうから、\textbf{電波}・\textbf{赤外線}・\textbf{可視光}（目に見える光）・\textbf{紫外線}・\textbf{X 線}・\textbf{γ 線}（ガンマ線）とよび分けています（図\ref{ph:電-15:fig-kazoku}）。\bkidx{でんぱ}{電波}どれも真空の中を同じ速さ $c$ で進む、同じなかまの波です。

\begin{bkfigure}
\begin{tikzpicture}[x=1mm, y=1mm]
  % 横軸：波長の対数。x = 7.5*(4 - log10(λ/m))  （λ=10^4 m で 0、10^-12 m で 120）
  \fill[bkBlue!25] (0,0) rectangle (60,10);
  \fill[bkRed!30] (60,0) rectangle (75.8,10);
  \shade[left color=bkRed!70, right color=bkPurple!70] (75.8,0) rectangle (78.2,10);
  \fill[bkPurple!25] (78.2,0) rectangle (90,10);
  \fill[bkGray!35] (90,0) rectangle (108,10);
  \fill[bkGreen!30] (108,0) rectangle (120,10);
  \draw[bkInk, line width=0.6pt] (0,0) rectangle (120,10);
  \node[sxp lab] at (30,5) {電波};
  \node[sxp lab] at (67.9,5) {赤外線};
  \node[sxp lab] at (84.1,5) {紫外線};
  \node[sxp lab] at (99,5) {X 線};
  \node[sxp lab] at (114,5) {γ 線};
  \draw[bkInk, line width=0.5pt] (77,10) -- (77,13) node[anchor=south, sxp lab] {可視光（\qtyrange{380}{780}{nm}）};
  % 波長の目もり
  \foreach \e/\t in {4/{$10^4$},2/{$10^2$},0/{1},-2/{$10^{-2}$},-4/{$10^{-4}$},-6/{$10^{-6}$},-8/{$10^{-8}$},-10/{$10^{-10}$},-12/{$10^{-12}$}} {
    \draw[bkInk, line width=0.5pt] ({7.5*(4-\e)},0) -- ({7.5*(4-\e)},-1.5);
    \node[sxp lab, anchor=north] at ({7.5*(4-\e)},-2) {\t};}
  \node[sxp lab, anchor=north west] at (0,-6.5) {波長（m）：右ほど短い};
  % 使いみち
  \foreach \x/\t in {11/{AM\\ラジオ},26/{FM・\\テレビ},43/{スマホ・\\電子レンジ},66/{リモコン・\\こたつ},85/{日焼け},101/{レント\\ゲン}} {
    \node[sxp lab, anchor=north, align=center] at (\x,-11.5) {\t};}
\end{tikzpicture}
\bkfigcaption{電磁波の家族（横は波長。1目もりで100分の1になる対数の目もり：3-12）。境目ははっきり決まったものではなく、目安。電波は、法律（電波法）で周波数 \qty{3}{THz}（波長 \qty{0.1}{mm}）以下の電磁波と決められている。}
\label{ph:電-15:fig-kazoku}
\end{bkfigure}

\label{ph:電-15:clambda0}可視光は、波長約 \qty{380}{nm}（紫）から約 \qty{780}{nm}（赤）までのせまい範囲です（波-13 ①）。その外側の、目に見えない電磁波も、私たちの生活の中で使われています。

\begin{bkcenter}
\begin{bktable}{colspec={Q[l,m,wd=24mm] Q[c,m,wd=26mm] Q[l,m,wd=70mm]}, rowsep=1.2pt, row{1}={bg=bkTheme!10, font=\sffamily\small}, cells={font=\small}}
なかま & 波長の目安 & 身近な例・使いみち \\
電波 & \qty{0.1}{mm} より長い & ラジオ・テレビ・スマホ・Wi-Fi・電子レンジ・IC カード \\
赤外線 & 約 \qty{780}{nm}〜\qty{0.1}{mm} & テレビのリモコン、こたつ、体温をはかる放射温度計（熱-13） \\
可視光 & 約 \qtyrange{380}{780}{nm} & 目で見る光、光ファイバ（波-14 では赤外線も） \\
紫外線 & 約 \qtyrange{10}{380}{nm} & 日焼け、殺菌の灯 \\
X 線・γ 線 & 約 \qty{10}{nm} より短い & レントゲン写真、がんの治療（原-13） \\
\end{bktable}
\bktabcaption{電磁波の家族と使いみち。X 線と γ 線は、波長ではなく、出るしくみ（電子か原子核か）で分けることが多い（原-12）。}
\label{ph:電-15:tab-kazoku}
\end{bkcenter}

\bksubsection{2}{身の回りの電波の周波数と波長}

\leavevmode\label{ph:電-15:sec2}電波は、使う周波数が法律（電波法）で決められていて、同じ周波数をほかの人が勝手に使うことはできません。身の回りの電波の周波数と、$c=f\lambda$ で計算した波長を並べます（光の速さ $c=3.00\times10^8$\,m/s）。

\begin{bkcenter}
\begin{bktable}{colspec={Q[l,m,wd=40mm] Q[c,m,wd=36mm] Q[c,m,wd=40mm]}, rowsep=1.2pt, row{1}={bg=bkTheme!10, font=\sffamily\small}, cells={font=\small}}
電波 & 周波数 $f$ & 波長 $\lambda$ \\
IH 調理器 & \qtyrange{20.05}{100}{kHz} & \qty{15}{km}〜\qty{3}{km} \\
AM（中波）ラジオ & \qtyrange{526.5}{1606.5}{kHz} & 約 \qtyrange{570}{190}{m} \\
交通系 IC カード & \qty{13.56}{MHz} & 約 \qty{22}{m} \\
FM ラジオ & \qtyrange{76}{95}{MHz} & 約 \qtyrange{3.9}{3.2}{m} \\
地上デジタルテレビ & \qtyrange{470}{710}{MHz} & 約 \qtyrange{64}{42}{cm} \\
Wi-Fi & 約 \qty{2.4}{GHz}・\qty{5}{GHz} & 約 \qty{12.5}{cm}・\qty{6}{cm} \\
電子レンジ & \qty{2450}{MHz}（±\qty{50}{MHz}） & 約 \qty{12.2}{cm} \\
\end{bktable}
\bktabcaption{身の回りの電波（周波数は電波法施行規則・総務省の資料。IH 調理器は電波を出すためではなく、電-8 ⑫ のうず電流で温める道具だが、周波数は法律で決められている）。}
\label{ph:電-15:tab-denpa}
\end{bkcenter}

\begin{phmijika}[電波の名前は、波長の桁で決まる]
AM ラジオの電波は数百 m、FM は数 m、スマホや電子レンジは十数 cm です。工学では、波長の桁で、長波（km）・中波（hm）・短波（10 m）・超短波（m）・極超短波（dm）・マイクロ波（cm〜mm）などとよび分けます。ラジオの「中波」「短波」は、この名前です。波長が長いほど、建物や山のかげに回りこみやすく（波-8 ④の回折）、短いほど、たくさんの情報を運べます（波-14 ⑧）。
\end{phmijika}

\phstage{te}{ゆれ方とエネルギーを計算する}{3}

\bksubsection{3}{電磁波の速さ：$c=f\lambda$}

\leavevmode\label{ph:電-15:sec3}電磁波は、真空の中をどれも同じ速さ $c$ で進みます。波の速さ・周波数・波長の関係は、ほかの波と同じ（波-5 ②）で
\pheq{c=f\lambda}{m/s ＝ Hz × m}
\label{ph:電-15:clambda}です。真空の中の光の速さは、今では
\[ c=\qty{299792458}{m/s}\qquad（\text{約}\ 3.00\times10^8\ \mathrm{m/s}） \]
と、はかる値ではなく\textbf{決めた値}になっています（1983年から。1 m の長さを、光が $1/299792458$ 秒に進む長さとして決めた：波-13 ⑦）。1秒に地球を7周半する速さです。

\begin{phyomi}[$c=f\lambda$]
「1秒に $f$ 回ゆれ、1回のゆれで波長 $\lambda$ だけ進むので、1秒に進む長さは $f\lambda$」。電磁波では $c$ が決まっているので、周波数と波長は反比例します。周波数を2倍にすると、波長は半分です（例題1・2）。物の中では速さが $c/n$（$n$ は屈折率：波-12 ③）になり、周波数は変わらないので、波長が $1/n$ になります。
\end{phyomi}

\bksubsection{4}{電場と磁場が直角にゆれる横波：$E=cB$}

\leavevmode\label{ph:電-15:sec4}電磁波の中では、電場と磁場がどんな向きにゆれているのでしょう。電-14 の積分形の法則だけで調べます。

いちばん簡単な場合として、図\ref{ph:電-15:fig-step}のように、ある瞬間にスイッチを入れて、境目（\textbf{波面}）の後ろだけに、一様な電場 $\vec E$（上向き）と一様な磁場 $\vec B$（紙面の手前向き）がある場が、右へ速さ $v$ で進んでいるとします。波面の前には、まだ何もありません。この場が、マクスウェル方程式をみたすための条件を探します。

\begin{bkfigure}
\begin{tikzpicture}[x=1mm, y=1mm]
  \fill[bkTheme!10] (0,0) rectangle (60,34);
  \draw[bkInk, line width=1.0pt] (60,0) -- (60,34);
  \node[sxp lab, anchor=south] at (60,34.5) {波面};
  \draw[sxp ray, line width=1.2pt] (61,30) -- (72,30); \node[sxp lab, anchor=west] at (72.5,30) {速さ $v$};
  \foreach \x in {6,18} {\draw[sxp E] (\x,6) -- (\x,28);}
  \node[sxp lab, text=bkRed!80!black, anchor=west] at (6.6,30) {$\vec E$};
  \foreach \x in {12,24,36,48} {\foreach \y in {11.5,17,22.5} {\node[text=bkBlue!80!black, font=\footnotesize] at (\x,\y) {$\odot$};}}
  \node[sxp lab, text=bkBlue!80!black, anchor=north west] at (24,6) {$\vec B$（手前向き）};
  % 長方形の道
  \draw[bkPurple!80!black, line width=1.1pt] (40,8) rectangle (80,26);
  \draw[bkPurple!80!black, line width=1.1pt, -{Stealth[length=3pt]}] (60.5,26) -- (64,26);
  \node[sxp lab, text=bkPurple!80!black, anchor=west] at (80.6,17) {道（高さ $h$）};
  \draw[sxp dim] (60,-2) -- (70,-2); \node[sxp lab, anchor=west] at (70.6,-2) {1秒で $v$ 進む};
  \node[sxp lab, anchor=north west, align=left] at (88,32) {後ろ：$\vec E$ と $\vec B$ がある\\前：まだ何もない};
\end{tikzpicture}
\bkfigcaption{右へ進む電磁場の波面（1つの瞬間）。電場（実線）は上向き、磁場（$\odot$ は紙面の手前向き）。波面をまたぐ長方形の道で、ファラデーの法則を使う。}
\label{ph:電-15:fig-step}
\end{bkfigure}

\Needspace{5\baselineskip}\noindent\textbf{ファラデーの法則から}\par\nobreak
図の長方形の道（高さ $h$）に、ファラデーの法則（電-14 ⑤の（3））を使います。道を1周する電場は、波面の後ろの、たての辺（長さ $h$）だけが $E\times h$ を作ります（前のたての辺には場がなく、横の辺は電場と直角）。一方、道を通る磁束は、波面が1秒に $v$ ずつ前へ進むので、1秒に $B\times(hv)$ ずつ増えます。2つが等しいので
\[ Eh=Bhv\qquad\text{より}\qquad E=vB \]

\Needspace{5\baselineskip}\noindent\textbf{アンペール・マクスウェルの法則から}\par\nobreak
こんどは、紙面に直角な面に、同じように波面をまたぐ長方形の道をとり、アンペール・マクスウェルの法則（電-14 ③）を使います。電流はないので、道を1周する磁場 $B\times h$ は、道を通る電場の流束の増える速さ $E\times(hv)$ に $\mu_0\varepsilon_0$ をかけたものです。
\[ Bh=\mu_0\varepsilon_0Ehv\qquad\text{より}\qquad B=\mu_0\varepsilon_0vE \]

\Needspace{5\baselineskip}\noindent\textbf{2つを合わせると、速さが決まる}\par\nobreak
1つめの式の $E=vB$ を2つめに入れると、$B=\mu_0\varepsilon_0v^2B$ です。両辺を $B$ でわると
\pheq{v^2=\frac{1}{\mu_0\varepsilon_0}\qquad\text{より}\qquad v=\frac{1}{\sqrt{\mu_0\varepsilon_0}}=c}{m/s ＝ 1/√((N/A²)(F/m))}
\label{ph:電-15:c}となりました。$\varepsilon_0=8.85\times10^{-12}$\,F/m、$\mu_0=4\pi\times10^{-7}$\,N/A² を入れると $3.00\times10^8$\,m/s で、光の速さです（8-14⑤）。電場と磁場の場は、勝手な速さでは進めず、ちょうど光の速さでだけ、つじつまが合うのです。そして
\pheq{E=cB}{V/m ＝ (m/s) × T}
\label{ph:電-15:EcB}です。\label{ph:電-15:yokonami}電場・磁場・進む向きは、たがいに直角です（図\ref{ph:電-15:fig-yoko}）。電磁波は\textbf{横波}です。\bkidx{よこなみ}{横波}波-13 ⑥で「光の電場は、進む向きに垂直な面の中でゆれる」と結果だけ使ったのは、このことです。電場から磁場へ右ねじを回すと、ねじは進む向きに進みます（$\vec S=\vec E\times\vec H$ が進む向き）。

\begin{bkfigure}
\begin{tikzpicture}[x=1mm, y=1mm]
  \draw[bkInk!60, line width=0.6pt, -{Stealth[length=2mm]}] (0,0) -- (112,0) node[right, sxp lab] {進む向き};
  \draw[bkRed!80!black, line width=1.3pt, domain=0:100, samples=80] plot (\x,{14*sin(\x*7.2)});
  \foreach \xx in {6.25,12.5,...,93.75} {\pgfmathsetmacro\sxpE{14*sin(\xx*7.2)}\draw[bkRed!80!black, line width=0.5pt] (\xx,0) -- (\xx,\sxpE);}
  \draw[bkBlue!80!black, line width=1.3pt, dash pattern=on 3pt off 1.5pt, domain=0:100, samples=80] plot ({\x+5.6*sin(\x*7.2)},{-4.2*sin(\x*7.2)});
  \foreach \xx in {6.25,12.5,...,93.75} {\pgfmathsetmacro\sxpB{sin(\xx*7.2)}\draw[bkBlue!80!black, line width=0.5pt] (\xx,0) -- ({\xx+5.6*\sxpB},{-4.2*\sxpB});}
  \node[sxp lab, anchor=west, text=bkRed!80!black] at (68,17.5) {電場 $\vec E$（実線。たて向きにゆれる）};
  \node[sxp lab, anchor=north, text=bkBlue!80!black] at (50,-15.8) {磁場 $\vec B$（破線。紙面に直角な向きを、ななめから見た形）};
  \draw[sxp dim] (12.5,22) -- (62.5,22); \node[sxp lab, anchor=south] at (37.5,22.4) {波長 $\lambda$};
\end{tikzpicture}
\bkfigcaption{電磁波（正弦波）のようす。電場と磁場は同じ山と谷の位置で（同じ位相で）ゆれ、どちらも進む向きに直角。大きさは $E=cB$。}
\label{ph:電-15:fig-yoko}
\end{bkfigure}

\begin{phbaisuru}{電場の振幅 $E_0$}
磁場の振幅も2倍になります（$B_0=E_0/c$）。電磁波の電場と磁場は、いつも決まった割合でいっしょにゆれます。\qty{1}{V/m} の電場には $3.3\times10^{-9}$\,T の磁場が組になっていて、磁場は地球の磁場（約 $5\times10^{-5}$\,T）の1万分の1もありません。電磁波の中の磁場は、とても小さく見えます。けれども、エネルギーで比べると、電場と磁場は同じだけ運んでいます（⑤）。
\end{phbaisuru}

\bksubsection{5}{電磁波が運ぶエネルギーと運動量}

\leavevmode\label{ph:電-15:sec5}電磁波は、太陽の熱を地球に運び、電子レンジの食べ物を温めます。エネルギーを運んでいるのです。

$E=cB$ を電-14 ⑦のエネルギー密度に入れると、$1/(\mu_0c^2)=\varepsilon_0$ なので、磁場の部分は
\[ \frac{B^2}{2\mu_0}=\frac{E^2}{2\mu_0c^2}=\frac12\varepsilon_0E^2 \]
で、電場の部分とちょうど同じです。合わせて $u=\varepsilon_0E^2$。このエネルギーが速さ $c$ で進むので、\qty{1}{m^2} の窓を1秒に通るエネルギー（ポインティングベクトルの大きさ）は $S=uc=c\varepsilon_0E^2$ です。正弦波では $E^2$ の時間の平均が $\tfrac12E_0^2$（4-9⑧と同じ $\sin^2$ の平均）なので、電磁波の\textbf{強さ}（平均のエネルギーの流れ）は
\pheq{I=\frac12c\varepsilon_0E_0^2}{W/m² ＝ (m/s)(F/m)(V/m)²}
\label{ph:電-15:I}です。強さは振幅の2乗に比例します。ロープの波や音（波-7 ①）と同じです。

\Needspace{5\baselineskip}\noindent\textbf{光は押す：放射圧}\par\nobreak
電磁波が、物の中の電子に当たったとします。電場 $\vec E$ は、電子を電場の向きにゆすります。動いている電子には、磁場 $\vec B$ からローレンツ力 $q\vec v\times\vec B$（電-6 ③）がはたらきます。$\vec v$ は電場の向き、$\vec B$ はそれに直角なので、この力は\textbf{電磁波の進む向き}を向きます（図\ref{ph:電-15:fig-oshi}）。
\begin{bkfigure}
\begin{tikzpicture}[x=1mm, y=1mm]
  \draw[sxp ray, line width=1.2pt] (0,12) -- (24,12); \node[sxp lab, anchor=south] at (12,13) {電磁波が進む向き};
  \filldraw[fill=bkGray!25, draw=bkInk, line width=0.6pt] (30,0) rectangle (70,28);
  \node[sxp lab, anchor=north] at (50,-1) {物（電磁波を吸う）};
  \filldraw[fill=bkBlue!60, draw=bkInk, line width=0.4pt] (52,12) circle (1.6);
  \draw[sxp E] (36,4) -- (36,22); \node[sxp lab, text=bkRed!80!black, anchor=west] at (37,20) {$\vec E$};
  \node[text=bkBlue!80!black] at (44,25) {$\odot$}; \node[sxp lab, text=bkBlue!80!black, anchor=west] at (45.5,25) {$\vec B$};
  \draw[bkGreen!50!black, line width=1.0pt, -{Stealth[length=3pt]}] (52,14) -- (52,21); \node[sxp lab, anchor=west, text=bkGreen!40!black] at (53,19.5) {$\vec v$};
  \draw[bkPurple!80!black, line width=1.2pt, -{Stealth[length=3pt]}] (54,12) -- (66,12);
  \node[sxp lab, anchor=west, text=bkPurple!80!black, align=left] at (72,12) {$q\vec v\times\vec B$\\（進む向きに押す）};
\end{tikzpicture}
\bkfigcaption{電磁波が物を押すしくみ。$\vec v$ は、電場にゆすられた電荷の速さ（正の電荷としてかいた。負の電荷では、ゆれの向きも電荷の符号も逆なので、力は同じく進む向き）。磁場（$\odot$）は紙面の手前向き。}
\label{ph:電-15:fig-oshi}
\end{bkfigure}
電子が1秒にもらうエネルギー（仕事率）は「力×速さ」$=qEv$、押される力は $qvB=qvE/c$ です。比べると、押す力は「もらう仕事率÷$c$」です。物が強さ $I$ の電磁波を全部吸うと、面積 $A$ がもらう仕事率は $IA$ なので、押す力は $IA/c$、圧力（\textbf{放射圧}）は
\pheq{p=\frac{I}{c}\qquad（\text{全部吸うとき。全部はね返すときは}\ 2I/c）}{Pa ＝ (W/m²) ÷ (m/s)}
\label{ph:電-15:p}です。\bkidx{ほうしゃあつ}{放射圧}はね返すときは、来た波を止める分と、逆向きに送り出す分で2倍になります。電磁波は、エネルギーだけでなく、運動量も運んでいるのです。光が粒（光子）であると考えたとき、1つの光子の運動量が「エネルギー÷$c$」になることを、物理シリーズの『原子』（原-2）で使います。

\begin{phmijika}[光の帆で進む宇宙ヨット]
太陽の光の放射圧はとても小さく、地球の近くで \qty{1}{m^2} あたり $10^{-5}$\,N ほどです（例題4）。けれども宇宙では空気の抵抗がないので、大きくて軽い帆を広げれば、少しずつ速さを増やせます。2010年に打ち上げられた日本の IKAROS（イカロス）は、差し渡し \qty{20}{m} の正方形の、厚さ \qty{7.5}{\micro m} のうすい膜を宇宙で広げ、太陽の光の放射圧で進む「ソーラーセイル」を実際に試しました。
\end{phmijika}

\phkokomade

\Needspace{9\baselineskip}
\phstage{shiki}{4つの式から電磁波を導く}{5}

\bksubsection{6}{波動方程式：光が電磁波であること}

\leavevmode\label{ph:電-15:sec6}④では、波面が一様な特別な場で調べました。どんな形の電磁波でも同じになることを、微分形で確かめます。電荷も電流もない真空の中では、マクスウェル方程式（電-14 ⑤）の（3）の両辺の rot をとり、（4）を入れ、$\mathrm{rot}\,\mathrm{rot}\,\vec E=\mathrm{grad}\,\mathrm{div}\,\vec E-\nabla^2\vec E$（8-12⑧）と $\mathrm{div}\,\vec E=0$ を使うと
\pheq{\frac{\partial^2\vec E}{\partial t^2}=\frac{1}{\mu_0\varepsilon_0}\,\nabla^2\vec E}{(V/m)/s² ＝ (m²/s²)(V/m)/m²}
\label{ph:電-15:hadou}が出ます（計算は 8-14⑤）。磁場も同じ式をみたします。これは、弦の波（波-5 ④）と同じ形の\textbf{波動方程式}で、速さは $c=1/\sqrt{\mu_0\varepsilon_0}$ です。

\begin{phflow}[マクスウェル方程式から光へ]
\phflowrow{マクスウェル方程式\\（電-14 ⑤）}{電荷も電流も\\ない真空}{$\dfrac{\partial^2\vec E}{\partial t^2}=c^2\nabla^2\vec E$}
\phflowrow{$c=\dfrac{1}{\sqrt{\mu_0\varepsilon_0}}$}{$\varepsilon_0$・$\mu_0$ の\\はかった値}{$c=3.00\times10^8$\,m/s\\（光の速さ）}
\phflowrow{$v=\dfrac{1}{\sqrt{\varepsilon\mu}}$\\（物の中）}{磁石につかない\\物（$\mu_{\mathrm r}=1$）}{$n=\dfrac{c}{v}=\sqrt{\varepsilon_{\mathrm r}}$}
\end{phflow}

音の波は、空気がおし縮められて、もとにもどろうとする力で伝わりました（波-5 ⑤）。空気がなければ、音は伝わりません。電磁波では、変わる電場が磁場を作り、変わる磁場が電場を作るので、ゆれを伝える物（媒質）が要りません。だから、太陽の光は、ほとんど真空の宇宙を通って地球へ届くのです（波-5 ①の約束）。

\Needspace{5\baselineskip}\noindent\textbf{物の中の速さと、屈折率}\par\nobreak
ガラスや水のような物の中では、電-14 ⑥のとおり $\varepsilon_0$ を $\varepsilon=\varepsilon_{\mathrm r}\varepsilon_0$ に、$\mu_0$ を $\mu=\mu_{\mathrm r}\mu_0$ にかえるので、速さは $v=1/\sqrt{\varepsilon\mu}=c/\sqrt{\varepsilon_{\mathrm r}\mu_{\mathrm r}}$ です。屈折率は $n=c/v$（波-12 ③）なので
\pheq{n=\sqrt{\varepsilon_{\mathrm r}\mu_{\mathrm r}}}{単位なし}
\label{ph:電-15:n}です。ガラスや水のように磁石につかない物では $\mu_{\mathrm r}=1$ で、$n=\sqrt{\varepsilon_{\mathrm r}}$ です。屈折率は、物が電場にどれだけ分極しやすいかで決まっていたのです（⑨・例題10）。

\Needspace{5\baselineskip}\noindent\textbf{正弦波の形}\par\nobreak
$x$ の向きに進む正弦波は、$E_y=E_0\sin(\omega t-kx)$、$B_z=B_0\sin(\omega t-kx)$ と書けます（8-13⑩）。波数 $k=2\pi/\lambda$ と $\omega=ck$ の関係は、ほかの波と同じです（波-5 ③）。電場と磁場は同じ $\sin$ の中身をもち、いっしょに山と谷になります（図\ref{ph:電-15:fig-yoko}）。

\bksubsection{7}{平面波と、波のインピーダンス}

\leavevmode\label{ph:電-15:sec7}波-7 では、ロープの波の「力÷速さ」、音の「音圧÷速さ」を\textbf{インピーダンス}とよびました。「押すもの÷流れるもの」です。電磁波では、電場が「押すもの」（電圧のなかま）、磁場の強さ $H$ が「流れるもの」（電流のなかま）にあたります。その比を、\textbf{波のインピーダンス} $\eta$（イータ）といいます。\bkidx{なみのいんぴーだんす}{波のインピーダンス}④の $E=cB=c\mu_0H$ から、真空では
\pheq{\eta_0=\frac{E}{H}=\mu_0c=\sqrt{\frac{\mu_0}{\varepsilon_0}}\approx\qty{377}{\ohm}}{Ω ＝ (V/m) ÷ (A/m)}
\label{ph:電-15:eta}です。物の中では $\eta=\sqrt{\mu/\varepsilon}$ で、磁石につかない物では $\eta=\eta_0/\sqrt{\varepsilon_{\mathrm r}}=\eta_0/n$ です。波-7 ⑨で「光にとってのインピーダンスは屈折率に反比例する（$Z\propto1/n$）」と結果だけ使ったのは、このことです。

\begin{phyomi}[$\eta_0\approx\qty{377}{\ohm}$]
「真空の中の電磁波では、電場 \qty{1}{V/m} に、磁場の強さ $1/377$\,A/m がいつも組になっている」。単位が Ω なのは、V/m を A/m でわったからです。電線の上の電気の波の特性インピーダンス $Z_0$（電-13 ⑦。同軸ケーブルの \qty{50}{\ohm} など）が「電圧÷電流」なのと同じ考えで、空間そのものが、電磁波に対して 377 Ω の「線」になっているのです。強さは $I=\tfrac12E_0H_0=E_0^2/(2\eta_0)$ とも書けます（例題5）。
\end{phyomi}

\bksubsection{8}{ゆれる電荷が出す電磁波：ラーモアの式}

\leavevmode\label{ph:電-15:sec8}電磁波は、どうやって生まれるのでしょう。止まった電荷の周りには電場だけがあり、一定の速さで動く電荷の周りには電場と磁場がありますが、どちらも電荷といっしょに動くだけで、遠くへ離れていきません。電磁波が生まれるのは、電荷が\textbf{加速度運動}をするときです。波-13 ⑧で結果だけ使った「出す光の強さは加速度の2乗に比例する」を、ここで導きます。

\Needspace{5\baselineskip}\noindent\textbf{電気力線の折れ曲がり}\par\nobreak
止まっていた電荷 $q$ が、時刻0から短い時間 $\Delta t$ だけ加速度 $a$ で動きはじめ、そのあとは小さな速さ $v=a\,\Delta t$（$v\ll c$）で進むとします。電場の変化は、光の速さ $c$ でしか伝わりません。そのため、時刻 $t$ に見ると（図\ref{ph:電-15:fig-ore}）
\begin{itemize}[leftmargin=1.4\zw, itemsep=1pt, topsep=2pt, label={\color{bkTheme}●}]
\item 半径 $ct$ より外：まだ「電荷が動いた」知らせが届いていない。電気力線は、はじめの位置 O からまっすぐ出ている。
\item 半径 $c(t-\Delta t)$ より内（加速が終わった知らせが届いた所）：一定の速さで動く電荷の電場で、電気力線は今の位置 P からまっすぐ出ている（$v\ll c$ のとき）。
\item その間の、厚さ $c\,\Delta t$ のうすい殻：外と内の電気力線を、つなぐ部分。電気力線はとぎれない（電-2 ③のガウスの法則）ので、ここで\textbf{折れ曲がる}。
\end{itemize}

\begin{bkfigure}
\begin{tikzpicture}[x=1mm, y=1mm]
  \def\R{26}\def\r{22}
  \draw[bkGray, line width=0.4pt, dash pattern=on 2pt off 1.5pt] (0,0) circle (\R);
  \draw[bkGray, line width=0.4pt, dash pattern=on 2pt off 1.5pt] (0,0) circle (\r);
  % 外の線（O から）と内の線（P から。内の点線の円まで）
  \foreach \a in {30,60,120,150,210,240,300,330} {
    \pgfmathsetmacro\sxps{-5*cos(\a)+sqrt(484-25*(sin(\a))^2)}
    \draw[bkRed!80!black, line width=0.8pt] (\a:\R) -- (\a:32);
    \draw[bkRed!80!black, line width=0.8pt] (5,0) -- ($(5,0)+(\a:\sxps)$) -- (\a:\R);}
  \fill[bkInk] (0,0) circle (0.8); \node[sxp lab, anchor=east] at (-1.2,0) {O};
  \filldraw[fill=bkRed!70, draw=bkRed!80!black] (5,0) circle (1.4); \node[sxp lab, anchor=north] at (5,-1.8) {P};
  \node[sxp lab, anchor=north, align=center] at (0,-33) {外の点線：半径 $ct$\\内の点線：半径 $c(t-\Delta t)$};
  % 拡大図
  \begin{scope}[shift={(40,-14)}]
    \draw[bkInk, line width=0.5pt] (0,0) rectangle (38,30);
    \node[sxp lab, anchor=south] at (19,30.5) {折れ曲がりの拡大};
    \draw[bkRed!80!black, line width=1.0pt] (4,3) -- (12,12) -- (26,16) -- (33,27);
    \draw[sxp dim] (12,9) -- (26,9); \node[sxp lab, anchor=north] at (21,8.4) {横 $vt\sin\theta$};
    \draw[sxp dim] (29,16) -- (29,12.5); \node[sxp lab, anchor=west] at (29.6,13.5) {$c\,\Delta t$};
    \node[sxp lab, anchor=west, text=bkRed!80!black] at (2,24) {$E_\perp$ が大きい};
  \end{scope}
  \node[sxp lab, anchor=west, align=left] at (81,12) {外：O から};
  \node[sxp lab, anchor=west, align=left] at (81,4) {内：P から};
  \node[sxp lab, anchor=west, align=left] at (81,-4) {間の殻で\\折れ曲がる};
\end{tikzpicture}
\bkfigcaption{加速した電荷の周りの電気力線（動きを大きくかいた）。電荷は O から P へ右に動いた。外の線は O から、内の線は P から出ていて、その間の殻で折れ曲がる。折れ曲がった所では、進む向きに直角な電場 $E_\perp$ が強い。これが電磁波。}
\label{ph:電-15:fig-ore}
\end{bkfigure}

\Needspace{5\baselineskip}\noindent\textbf{折れ曲がった所の電場}\par\nobreak
動く向きから角 $\theta$ の向きの電気力線を見ます。殻の中で、電気力線は、半径の向きに $c\,\Delta t$ 進むあいだに、横へ $vt\sin\theta$ ずれます（O と P の距離 $vt$ の、横の向きの成分）。電場の向きは電気力線の向きなので、殻の中の電場の「横の成分 $E_\perp$」と「半径の成分 $E_r$」の比は、このずれの比です。
\[ \frac{E_\perp}{E_r}=\frac{vt\sin\theta}{c\,\Delta t}=\frac{a\,\Delta t\times t\sin\theta}{c\,\Delta t}=\frac{at\sin\theta}{c} \]
半径の成分は、止まった電荷と同じ $E_r=q/(4\pi\varepsilon_0r^2)$ です（殻の外と同じ本数の電気力線が通るから）。殻までの距離は $r=ct$ なので $t=r/c$ を入れて
\pheq{E_\perp=\frac{q}{4\pi\varepsilon_0r^2}\times\frac{ar\sin\theta}{c^2}=\frac{qa\sin\theta}{4\pi\varepsilon_0c^2r}}{V/m ＝ C(m/s²) ÷ ((F/m)(m/s)²m)}
\label{ph:電-15:Eperp}となります。ふつうの電場は $1/r^2$ で弱くなりますが、この横の電場は $1/r$ でしか弱くなりません。これが、加速度運動する電荷から遠くまで届く電磁波です。

\Needspace{5\baselineskip}\noindent\textbf{出す電磁波の仕事率}\par\nobreak
⑤で見たように、電磁波が運ぶエネルギーの流れは $S=c\varepsilon_0E_\perp^2$ です。半径 $r$ の球面の全体でたし合わせます。球面のタイルの面積は $r^2\sin\theta\,\dd\theta\,\dd\varphi$（8-8⑨）なので、$r^2$ が約分されて、
\[ P=\iint c\varepsilon_0E_\perp^2\,\dd A \]
\[ =c\varepsilon_0\Bigl(\frac{qa}{4\pi\varepsilon_0c^2}\Bigr)^2\times\int_0^{2\pi}\!\!\int_0^{\pi}\sin^3\theta\,\dd\theta\,\dd\varphi \]
です。後ろの積分は、$\cos\theta=u$ とおくと
\[ \int_0^\pi\sin^3\theta\,\dd\theta=\int_{-1}^{1}(1-u^2)\,\dd u=\frac43 \]
なので、$2\pi\times\tfrac43=\tfrac{8\pi}{3}$ です。まとめると
\pheq{P=\frac{q^2a^2}{6\pi\varepsilon_0c^3}}{W ＝ C²(m/s²)² ÷ ((F/m)(m/s)³)}
\label{ph:電-15:larmor}となります。これを\textbf{ラーモアの式}といいます。\bkidx{らーもあのしき}{ラーモアの式}1897年に、イギリスのラーモアが示しました。

\begin{phyomi}[$P=q^2a^2/(6\pi\varepsilon_0c^3)$]
「電荷が出す電磁波の仕事率は、加速度の2乗に比例する」。\label{ph:電-15:hassei}速さが一定なら（$a=0$）、電磁波は出ません。また、$E_\perp\propto\sin\theta$ なので、強さは $\sin^2\theta$ に比例し、\textbf{ゆれの向き（$\theta=0$）には出さず}、ゆれに直角な向きにいちばん強く出します（図\ref{ph:電-15:fig-sin2}）。電荷の2乗に比例するので、正の電荷でも負の電荷でも同じです。
\end{phyomi}

\begin{bkfigure}
\begin{tikzpicture}[x=1mm, y=1mm]
  \draw[bkGray, line width=0.4pt, -{Stealth[length=2pt]}] (-28,0) -- (28,0);
  \draw[bkGray, line width=0.4pt, -{Stealth[length=2pt]}] (0,-22) -- (0,22) node[anchor=south, sxp lab] {ゆれの向き};
  \draw[bkBlue!80!black, line width=1.2pt, domain=0:360, samples=120, smooth] plot ({20*(sin(\x))^2*sin(\x)},{20*(sin(\x))^2*cos(\x)});
  \draw[bkOrange!85!black, line width=1.2pt, {Stealth[length=2.4pt]}-{Stealth[length=2.4pt]}] (0,-4) -- (0,4);
  \filldraw[fill=bkRed!70, draw=bkRed!80!black] (0,0) circle (1.2);
  \draw[sxp ray] (0,0) -- (40:16); \node[sxp lab, anchor=south west] at (40:16) {$\theta$ の向き};
  \draw[bkInk, line width=0.4pt] (0,7) arc (90:40:7); \node[sxp lab] at (65:9.5) {$\theta$};
  \node[sxp lab, anchor=west, align=left] at (32,6) {原点からの距離が、その向きに\\出す強さ（$\sin^2\theta$ に比例）};
  \node[sxp lab, anchor=west, align=left] at (32,-8) {ゆれの向き（上下）には\\まったく出さない};
\end{tikzpicture}
\bkfigcaption{上下にゆれる電荷（またはアンテナ）が出す電磁波の、向きによる強さ。横から見ると8の字（立体ではドーナツの形）になる。}
\label{ph:電-15:fig-sin2}
\end{bkfigure}

\Needspace{5\baselineskip}\noindent\textbf{ゆれる電子と、空の青さ}\par\nobreak
光の電場にゆすられた電子は、$x=x_0\cos\omega t$ のように単振動します（波-13 ⑧の強制振動。振幅を、波-13 では $A$ と書いた）。加速度は $-\omega^2x$ で、振幅は $\omega^2x_0$ です（波-1 ②）。ラーモアの式から、出す光の仕事率は $(\omega^2x_0)^2=\omega^4x_0^2$ に比例します。波-13 ⑧で、可視光では $x_0$ が色によらずほぼ同じだったので、散乱の強さは $\omega^4$、つまり $1/\lambda^4$ に比例します。青（\qty{450}{nm}）は赤（\qty{700}{nm}）の約5.9倍強く散らされ、空は青く見えます（例題7）。波-13 で結果として使った「加速度の2乗」が、電場の折れ曲がりから導けました。

\begin{phmijika}[青空は偏光している：ブルースター角のわけ]
太陽の光で横向きにゆすられた空気の中の電子は、ゆれの向きには光を出しません。だから、太陽から $90°$ の向きの空の光は、ある1つの向きにゆれる光（偏光）が多くなります（やってみようの③）。波-13 ⑥のブルースター角も同じです。ガラスの表面近くの電子は、屈折光の電場の向きにゆれます。反射光が屈折光と直角になる入射角では、紙面の中でゆれる電子の\textbf{ゆれの向き}が、ちょうど反射光の進む向きになるので、その光は反射しません（⑨）。
\end{phmijika}

\begin{phmada}[宇宙から来る、一瞬の強い電波]
2007年、オーストラリアの電波望遠鏡の古い記録の中から、1000分の1秒ほどしか続かない、とても強い電波のパルスが見つかりました。その後、同じような電波がたくさん見つかり、\textbf{高速電波バースト}とよばれています。\bkidx{こうそくでんぱばーすと}{高速電波バースト}電波が宇宙の電子を通ってくるとき、周波数の低い電波ほど少しおくれて届くので（⑨の分散）、そのおくれから、多くは遠い銀河から来ていることがわかりました。

2020年には、私たちの銀河系の中の、とても強い磁場をもつ中性子星（マグネター）から、同じような電波が出たことが観測されました。けれども、遠い銀河から来る高速電波バーストが、みな同じしくみで出ているのか、何が電子をそれほど強く、そろって加速させるのかは、まだはっきりわかっていません。
\end{phmada}

\Needspace{10\baselineskip}\bksubsection{9}{物の中の電磁波と、表面での反射}

\leavevmode\label{ph:電-15:sec9}\Needspace{5\baselineskip}\noindent\textbf{水の比誘電率と屈折率}\par\nobreak
⑥で $n=\sqrt{\varepsilon_{\mathrm r}}$ を導きました。ところが水では、コンデンサーの実験ではかった比誘電率は約80（\qty{20}{\celsius}）なのに、$\sqrt{80}\approx8.9$ は、光の屈折率1.333（波-12 ③）とまるで合いません。これは、比誘電率が電場のゆれの速さ（周波数）で変わるからです。水の分子は、電場の向きにくるりと回って分極します（電-3 ⑥）。ゆっくり変わる電場なら回ってついていけますが、1秒に $10^{14}$ 回以上ゆれる光の電場には、分子の回転はとても追いつけません。光に対しては、分子の中の電子のずれだけが分極を作るので、比誘電率は $1.333^2\approx1.78$ と小さくなります（例題10）。周波数で屈折率が変わることが\textbf{分散}で、虹ができるわけでした（波-12 ⑧）。

\Needspace{5\baselineskip}\noindent\textbf{境目での電場と磁場のつながり方}\par\nobreak
光がガラスの表面に当たると、一部ははね返ります。その割合は、境目で電場と磁場がどうつながるかで決まります。電-3 ⑨・電-7 ⑦ と同じように、境目をはさむ細長い長方形の道に、ファラデーの法則とアンペール・マクスウェルの法則を使います。道の短い辺を0に近づけると、道を通る磁束・電場の流束は0になるので、\textbf{境目の面にそった電場の成分}と、\textbf{面にそった磁場の強さの成分}は、境目の両側で等しくなります（表面に電流が流れないとき）。

\Needspace{5\baselineskip}\noindent\textbf{真上から当たるとき}\par\nobreak
屈折率 $n_1$ の物から $n_2$ の物へ、境目に垂直に電磁波が当たるとします。来る波・はね返る波・通りぬける波の電場を $E_{\text{来}}$・$E_{\text{返}}$・$E_{\text{通}}$ とします（同じ向きを正）。電場は面にそっているので、つながり方の1つめから
\[ E_{\text{来}}+E_{\text{返}}=E_{\text{通}} \]
磁場の強さは $H=E/\eta=nE/\eta_0$（⑦）です。はね返る波は逆向きに進むので、$\vec E\times\vec H$ が逆を向くように、磁場の向きが逆になります。だから、つながり方の2つめは
\[ n_1(E_{\text{来}}-E_{\text{返}})=n_2E_{\text{通}} \]
\begin{bkfigure}
\begin{tikzpicture}[x=1mm, y=1mm]
  \fill[bkBlue!10] (60,0) rectangle (112,34);
  \draw[bkInk, line width=0.9pt] (60,0) -- (60,34);
  \node[sxp lab, anchor=north west] at (2,33.5) {屈折率 $n_1$（空気）};
  \node[sxp lab, anchor=north east] at (111,33.5) {屈折率 $n_2$（ガラス）};
  % 来る波
  \draw[sxp ray] (14,24) -- (34,24); \node[sxp lab, anchor=south] at (24,24.6) {来る波};
  \draw[sxp E] (42,19) -- (42,29); \node[sxp lab, text=bkRed!80!black, anchor=west] at (42.6,28) {$E_{\text{来}}$};
  \node[text=bkBlue!80!black] at (52,24) {$\odot$}; \node[sxp lab, text=bkBlue!80!black, anchor=south] at (52,25.6) {$H_{\text{来}}$};
  % はね返る波
  \draw[sxp ray] (34,9) -- (14,9); \node[sxp lab, anchor=north] at (24,8.4) {はね返る波};
  \draw[sxp E] (42,11) -- (42,7); \node[sxp lab, text=bkRed!80!black, anchor=west] at (42.6,8) {$E_{\text{返}}$};
  \node[text=bkBlue!80!black] at (52,14.5) {$\odot$}; \node[sxp lab, text=bkBlue!80!black, anchor=south] at (52,16) {$H_{\text{返}}$};
  % 通りぬける波
  \draw[sxp ray] (82,17) -- (102,17); \node[sxp lab, anchor=north] at (92,16.4) {通りぬける波};
  \draw[sxp E] (70,13) -- (70,21); \node[sxp lab, text=bkRed!80!black, anchor=west] at (70.6,20.5) {$E_{\text{通}}$};
  \node[text=bkBlue!80!black] at (77,27) {$\odot$}; \node[sxp lab, text=bkBlue!80!black, anchor=west] at (78.6,27) {$H_{\text{通}}$};
\end{tikzpicture}
\bkfigcaption{屈折率の小さい物から大きい物へ、垂直に当たる電磁波（$\odot$ は磁場の強さが紙面の手前向き）。境目では、電場どうし・磁場の強さどうしが、両側でつながる。はね返る波は、電場が逆向きになり、磁場の強さは来る波と同じ向きに加わる。}
\label{ph:電-15:fig-kyoukai}
\end{bkfigure}

2つから $E_{\text{通}}$ を消すと、電場で数えた反射係数は
\pheq{r_E=\frac{E_{\text{返}}}{E_{\text{来}}}=\frac{n_1-n_2}{n_1+n_2}=\frac{\eta_2-\eta_1}{\eta_1+\eta_2}}{単位なし}
\label{ph:電-15:rE}です。波-7 ④の「押すもので数えた反射係数」$r_{\text{押}}=(Z_2-Z_1)/(Z_1+Z_2)$ の $Z$ を $\eta$ にかえた式と、同じです。はね返るエネルギーの割合（\textbf{反射率}）は $R=r_E^{\,2}$ で、空気（$n_1=1$）からガラス（$n_2=1.5$）では $(0.5/2.5)^2=0.04$、つまり4\,\%です（例題8）。通りぬける割合は、強さが $nE^2$ に比例することから $n_2E_{\text{通}}^2/(n_1E_{\text{来}}^2)=1.5\times0.8^2=0.96$ で、合わせて1になります。

\Needspace{5\baselineskip}\noindent\textbf{山と谷が入れかわるわけ}\par\nobreak
\label{ph:電-15:irekawari}$n_2>n_1$（屈折率の大きい物との境で反射する）なら、$r_E$ は負です。はね返った光の電場は、来た光の電場と逆向き、つまり山と谷が入れかわります。$n_2<n_1$ なら $r_E$ は正で、そのままです。波-13 ④の表（シャボン膜の干渉で使った「反射で山と谷が入れかわるとき」）は、この式から出ていたのです。屈折率の大きい物では、同じ電場に対して磁場の強さが大きい（$H=nE/\eta_0$）ので、境目の「流れるもの」の側を合わせるために、はね返る波の電場が逆向きになるのです。波-7 ⑤で、軽いひもから重いひもへ進む波の変位が山から谷にかわったのと同じしくみです。

\Needspace{5\baselineskip}\noindent\textbf{ななめに当たるとき：フレネルの式}\par\nobreak
光がななめに当たるときも、同じつながり方から反射率が出ます。計算は長いので、結果だけ示します（\textbf{フレネルの式}）。\bkidx{ふれねるのしき}{フレネルの式}入射角を $\theta_1$、屈折角を $\theta_2$（$n_1\sin\theta_1=n_2\sin\theta_2$：波-12 ③）とすると、電場が紙面に垂直にゆれる光と、紙面の中でゆれる光の電場の反射係数は
\[ r_{\perp}=\frac{n_1\cos\theta_1-n_2\cos\theta_2}{n_1\cos\theta_1+n_2\cos\theta_2},\qquad r_{\parallel}=\frac{n_2\cos\theta_1-n_1\cos\theta_2}{n_2\cos\theta_1+n_1\cos\theta_2} \]
で、反射率はその2乗です。$\theta_1=0$ ではどちらも真上からの式にもどります（$r_\parallel$ は向きの決め方で符号が逆になりますが、2乗は同じ）。

\begin{bkfigure}
\begin{tikzpicture}
\begin{axis}[bkplot, width=108mm, height=44mm, xmin=0, xmax=90, ymin=0, ymax=1.0, xtick={0,15,30,45,56.3,75,90}, xticklabels={0,15,30,45,\raisebox{-2.2mm}{56.3},75,90}, ytick={0,0.5,1}, xlabel={入射角 $\theta_1$（°）}, ylabel={反射率}]
  \addplot[bkcurve, domain=0:89, samples=90] {((cos(x)-sqrt(2.25-sin(x)^2))/(cos(x)+sqrt(2.25-sin(x)^2)))^2};
  \addplot[bkcurve2, dashed, domain=0:89, samples=90] {((2.25*cos(x)-sqrt(2.25-sin(x)^2))/(2.25*cos(x)+sqrt(2.25-sin(x)^2)))^2};
  \node[font=\sffamily\footnotesize, anchor=west, align=left, fill=white, inner sep=1.5pt] at (axis cs:4,0.7) {実線：紙面に垂直なゆれ（$r_\perp^{\,2}$）\\破線：紙面の中のゆれ（$r_\parallel^{\,2}$）};
\end{axis}
\end{tikzpicture}
\bkfigcaption{空気からガラス（$n=1.5$）へ入る光の反射率（フレネルの式。波-13 の図と同じ）。真上から（$0°$）はどちらも4\,\%。紙面の中のゆれは、ブルースター角 $56.3°$ でちょうど0になる。}
\label{ph:電-15:fig-fresnel}
\end{bkfigure}

$r_\parallel=0$ になるのは、分子が0になる
\[ n_2\cos\theta_1=n_1\cos\theta_2 \]
のときです。これと屈折の法則 $n_1\sin\theta_1=n_2\sin\theta_2$ の、左辺どうし・右辺どうしをかけると
\[ n_1n_2\sin\theta_1\cos\theta_1=n_1n_2\sin\theta_2\cos\theta_2 \]
つまり
\[ \sin2\theta_1=\sin2\theta_2\qquad\text{から}\qquad \theta_1+\theta_2=90° \]
が出ます（$\theta_1\neq\theta_2$ のとき）。反射光と屈折光が直角になる入射角で、$\tan\theta_{\mathrm B}=n_2/n_1$（ブルースター角：波-13 ⑥）です。これは、⑧の「ゆれる電子はゆれの向きに光を出さない」と同じことを、場のつながり方から確かめたものです。

\bksubsection{10}{金属と電磁波：表皮深さと遮へい}

\leavevmode\label{ph:電-15:sec10}金属の中には、自由に動ける電子がたくさんあります。電磁波の電場が金属に入ろうとすると、電子が動いて電流が流れ、電場を打ち消してしまいます。電-8 ⑫で学んだとおり、電磁波は表面から\textbf{表皮深さ}
\pheq{\delta=\sqrt{\frac{2}{\omega\mu\sigma}}}{m ＝ √(1/((1/s)(N/A²)(S/m)))}
\label{ph:電-15:delta}くらいまでしか入れず、その間に急に弱まります（$\sigma$ は導電率）。銅では、\qty{50}{Hz} で約 \qty{9}{mm}、\qty{2.45}{GHz} では約 \qty{1.3}{\micro m} です（例題11）。

入れなかった電磁波は、どうなるのでしょう。金属の波のインピーダンスは、空気の 377 Ω よりずっと小さく（銅で \qty{2.45}{GHz} なら $0.02$ Ω ほど）、⑨の $r_E=(\eta_2-\eta_1)/(\eta_1+\eta_2)$ はほぼ $-1$ です。電磁波は、ほとんど全部はね返されます。金属の板や箱で、電磁波を通さないようにすることを\textbf{電磁遮へい}（シールド）といいます。\bkidx{でんじしゃへい}{電磁遮へい}アルミはくで包んだスマホがつながりにくくなるのは、このためです（やってみよう②）。

\Needspace{5\baselineskip}\noindent\textbf{穴のあいた金属の網}\par\nobreak
では、穴のあいた網はどうでしょう。金属の筒（穴）の中を電磁波が通るには、筒の中で電場が、金属の壁で0になる形のゆれ（定在波：波-6 ②）を作らなければなりません。直径 $d$ の円い穴では、波長が
\pheq{\lambda_{\mathrm c}\approx1.71\,d}{m ＝ m}
\label{ph:電-15:lc}（\textbf{遮断波長}）より長い電磁波は、筒の中に入るとすぐ弱まり、奥へ進めません（\num{1.71} はベッセル関数という特別な関数から出る数で、結果だけ使います）。\bkidx{しゃだんはちょう}{遮断波長}そのとき、穴の中で電場は、長さ $\lambda_{\mathrm c}/(2\pi)$ 進むごとにおよそ $1/e$（約0.37倍）に弱まります。

\begin{bkfigure}
\begin{tikzpicture}[x=1mm, y=1mm]
  % 網
  \filldraw[fill=bkGray!40, draw=bkInk, line width=0.6pt] (50,0) rectangle (54,40);
  \foreach \y in {4,10,...,36} {\fill[white] (50,\y-1) rectangle (54,\y+1);}
  \node[sxp lab, anchor=north] at (52,-1) {網（穴の直径 $d$）};
  % マイクロ波
  \draw[bkRed!80!black, line width=1.1pt, domain=0:48, samples=80] plot (\x,{20+10*sin(9*\x)});
  \draw[bkRed!80!black, line width=1.1pt, domain=0:48, samples=80, dash pattern=on 2pt off 1.5pt] plot (\x,{20+6*sin(9*\x+180)});
  \node[sxp lab, anchor=south west, text=bkRed!80!black] at (0,31) {マイクロ波（波長 \qty{12.2}{cm}）};
  \node[sxp lab, anchor=north west, text=bkRed!80!black, align=left] at (2,9) {ほとんどはね返る\\（破線）};
  \draw[bkRed!80!black, line width=0.9pt, domain=54:70, samples=30] plot (\x,{20+0.6*sin(9*\x)});
  \node[sxp lab, anchor=west, align=left] at (72,22) {外へはほとんど\\出ない};
  % 光
  \foreach \y in {4,10,16} {\draw[bkOrange!85!black, line width=0.8pt, -{Stealth[length=2.4pt]}] (40,\y) -- (64,\y);}
  \node[sxp lab, anchor=west, text=bkOrange!75!black, align=left] at (72,8) {光（波長 約 \qty{0.5}{\micro m}）\\は穴を通りぬける};
\end{tikzpicture}
\bkfigcaption{電子レンジの窓の網（大きさは目安。網の穴はとても小さくかいた）。穴の直径が波長よりずっと小さいので、マイクロ波は穴を通れずはね返る。光の波長は穴よりずっと小さいので、通りぬける。}
\label{ph:電-15:fig-ami}
\end{bkfigure}

\Needspace{5\baselineskip}\noindent\textbf{はてな？の答え}\par\nobreak
\label{ph:電-15:hatena}\bkhatenaref に答えます。電子レンジのマイクロ波は \qty{2.45}{GHz} で、波長は \qty{12.2}{cm} です（例題1）。窓の網の穴は直径が数 mm 以下で、遮断波長（直径の約1.7倍）は数 mm です。マイクロ波の波長は、それよりずっと長いので、穴の中を進めず、金属の網ではね返されます（直径 \qty{1}{mm}・厚さ \qty{1}{mm} の穴なら、通りぬける強さは約1600分の1：例題12）。一方、光の波長は \qty{0.5}{\micro m} ほどで、穴よりずっと小さいので、ふつうに通りぬけます。網は、穴の大きさで、通す電磁波と通さない電磁波を分ける「ふるい」なのです。

\bksubsection{11}{アンテナ：電流のゆれを電波にする}

\leavevmode\label{ph:電-15:sec11}アンテナは、⑧の「ゆれる電荷」をたくさん集めたものです。波-14 ⑤で見た半波長のアンテナでは、電線の中の電子が、電波の周波数で上下にゆれ、電流の定在波（まん中で大きく、はしで0）ができます。電子が加速度運動をするので、⑧のとおり電磁波を出します。強さは $\sin^2\theta$ に近い形で、アンテナの向き（電線の向き）にはほとんど出ず、電線に直角な向きに強く出ます（図\ref{ph:電-15:fig-sin2}）。テレビの受信アンテナや、基地局のアンテナを、電波の来る向きに向けて立てるのは、このためです。

アンテナから出ていく電力は、電流の実効値の2乗に比例します。その比例の定数を\textbf{放射抵抗} $R_{\mathrm r}$ といいます。\bkidx{ほうしゃていこう}{放射抵抗}
\pheq{P=R_{\mathrm r}\times（\text{根もとの電流の実効値}）^2}{W ＝ Ω × A²}
\label{ph:電-15:Rr}「抵抗」といっても、熱を出すのではなく、電力を電波にして空間へ送り出す分を、抵抗におきかえて表したものです。細い半波長のアンテナでは、マクスウェル方程式から計算すると $R_{\mathrm r}\approx\qty{73}{\ohm}$ です（例題13）。アンテナをつなぐケーブルの特性インピーダンス（\qty{50}{\ohm}・\qty{75}{\ohm}：電-13）と近い値なので、はね返りが少なく、電力をよく送りこめます（波-7 ⑧の整合）。

\begin{phmijika}[IC カードは、電波ではなく磁場で話す]
交通系 IC カードの \qty{13.56}{MHz} の波長は約 \qty{22}{m} で、カードの中のコイル（数 cm）は、波長よりずっと小さいアンテナです。小さいアンテナは電磁波をほとんど出しません。IC カードは、改札機のコイルの近くの、変わる磁場（電-8 ⑬の相互誘導）でやりとりしています。だから、カードを数 cm まで近づけないと反応しないのです。
\end{phmijika}

\bksubsection{12}{電磁ノイズと、電波の安全}

\leavevmode\label{ph:電-15:sec12}電流が速く変わる所では、⑧のとおり、意図しなくても電磁波が出ます。パソコンや、モーターを速く切りかえる回路（電-16 のインバーター）、静電気の放電（電-1）のパチッという一瞬の大きな電流は、広い周波数の電磁波（\textbf{電磁ノイズ}）を出し、ラジオに雑音を入れたり、近くの機械の動きを乱したりします。電-14 の工学の駅の EMC（電磁両立性）は、このノイズを出さないように、受けても乱れないように、配線や遮へい（⑩）をくふうする技術です。

電波を出す設備の周りで、人が受ける電波の強さには、法律で上限が決められています（電波防護指針。電波法施行規則 第21条の4）。たとえば \qty{1.5}{GHz} から \qty{300}{GHz} の電波では、6分間の平均で、ポインティングベクトルの大きさが \qty{1}{mW/cm^2}（$=\qty{10}{W/m^2}$）、電場が \qty{61.4}{V/m} です。$S=E^2/\eta_0$ で確かめると、$61.4^2\div377\approx\qty{10}{W/m^2}$ で、⑦の関係と合っています。この値は、体温が上がりはじめるなどの影響が出る強さに、約50倍の余裕（安全率）をとって決められています（例題14）。

%% ---------------- たいせつ・ことば・記号 ----------------
\begin{taisetsu}
\begin{itemize}[leftmargin=1.2\zw, itemsep=2pt, topsep=0pt, label={\color{bkNavy}●}]\raggedright
\item 電波・赤外線・可視光・紫外線・X 線・γ 線は、波長の違う同じ電磁波。真空の中の速さ $c=\qty{299792458}{m/s}$（決めた値）、$c=f\lambda$。
\item 電磁波は横波。電場・磁場・進む向きはたがいに直角で、$E=cB$。速さは $c=1/\sqrt{\mu_0\varepsilon_0}$。
\item 強さ $I=\tfrac12c\varepsilon_0E_0^2=E_0^2/(2\eta_0)$。放射圧 $p=I/c$（はね返すと $2I/c$）。
\item 物の中では $n=\sqrt{\varepsilon_{\mathrm r}\mu_{\mathrm r}}$、波のインピーダンス $\eta=\eta_0/n$（真空 $\eta_0\approx\qty{377}{\ohm}$）。
\item 加速度運動する電荷が電磁波を出す。仕事率 $P=q^2a^2/(6\pi\varepsilon_0c^3)$（加速度の2乗）。ゆれの向きには出さない（$\sin^2\theta$）。
\item 真上から当たるときの電場の反射係数 $r_E=(n_1-n_2)/(n_1+n_2)$。屈折率の大きい物での反射で山と谷が入れかわる。ブルースター角 $\tan\theta_{\mathrm B}=n_2/n_1$。
\item 金属は表皮深さ $\delta$ より深くは電磁波を入れず、ほとんどはね返す。直径 $d$ の穴は、波長が約 $1.71d$ より長い電磁波を通さない。
\end{itemize}
\end{taisetsu}

\begin{kotoba}\small
\bkword{電磁波}{電場と磁場のゆれが、たがいを作りながら真空の中を速さ $c$ で進む横波。}
\bkword{電波}{周波数 \qty{3}{THz} 以下の電磁波（電波法の定義）。}
\bkword{波のインピーダンス}{電磁波の電場÷磁場の強さ。真空で約 \qty{377}{\ohm}。}
\bkword{放射圧}{電磁波が物を押す圧力。全部吸うと $I/c$。}
\bkword{ラーモアの式}{加速度運動する電荷が出す電磁波の仕事率。加速度の2乗に比例。}
\bkword{フレネルの式}{境目での電磁波の反射係数・透過係数を、入射角と屈折率で表す式。}
\bkword{電磁遮へい}{金属で囲んで、電磁波を通さないようにすること。}
\bkword{遮断波長}{金属の管や穴を、電磁波が通れるいちばん長い波長。}
\bkword{放射抵抗}{アンテナが電波にして出す電力を、抵抗におきかえた値。}
\bkword{高速電波バースト}{宇宙から届く、1000分の1秒ほどのとても強い電波のパルス。}
\end{kotoba}
\bkkigocards{\bkkigocard{$\eta$}{イータ}{波のインピーダンス〔Ω〕。電場÷磁場の強さ。真空で $\eta_0\approx\qty{377}{\ohm}$。電-13 の特性インピーダンス $Z_0$、電-12 の効率 $\eta$ とは別の量}{wave impedance}{電-15}}{\bkkigocard{$\lambda_{\mathrm c}$}{ラムダ・シー}{遮断波長〔m〕。金属の管や穴を通れるいちばん長い波長}{cutoff wavelength}{電-15}}

\begin{senmonsho}[平面波と放射]
\begin{bkhonbun}
無損失媒質中の平面波では $\symbfit{E}$，$\symbfit{H}$，伝搬方向が右手系をなし，$|\symbfit{E}|/|\symbfit{H}|=\eta=\sqrt{\mu/\varepsilon}$ を波動インピーダンス（自由空間では $\eta_0=120\pi\simeq377\,\Omega$）という。加速度 $\symbfit{a}$ の点電荷の全放射電力はラーモアの公式 $P=q^2a^2/(6\pi\varepsilon_0c^3)$ で与えられ，微小ダイポールの放射パターンは $\sin^2\theta$ に比例する。
\end{bkhonbun}
\bkyaku{無損失媒質＝エネルギーを吸わない物。伝搬方向＝進む向き。「右手系をなす」は、$\vec E$ から $\vec H$ へ右ねじを回すと進む向きに進むこと。$120\pi$ は $\mu_0c$ を $\mu_0=4\pi\times10^{-7}$、$c=3\times10^8$ で計算した値（$\approx377$）。「全放射電力」はこの本の仕事率 $P$。「微小ダイポール」は、ゆれる電荷・短いアンテナ。「放射パターン」は、向きによる強さの形（図\ref{ph:電-15:fig-sin2}）。}
\end{senmonsho}

%% ---------------- 例題 ----------------
\bksection{reidai}{例題でたしかめよう}

\noindent 定数は $c=3.00\times10^8$\,m/s、$\varepsilon_0=8.85\times10^{-12}$\,F/m、$\mu_0=4\pi\times10^{-7}$\,N/A²、$\eta_0=\qty{377}{\ohm}$、電気素量 $e=1.60\times10^{-19}$\,C とします。

\begin{phreidai}[電子レンジのマイクロ波の波長]
\label{ph:電-15:ex1}電子レンジは、周波数 \qty{2.45}{GHz} の電磁波で食べ物を温めます。波長を求めましょう。電子レンジの中にできる定在波の、腹と腹の間はいくらですか。
\phbamen
はてな？の電子レンジです（波-6 ⑧ の例題11と同じ数）。
\phmushi
レンジの中は空気で、空気の屈折率（1.0003：波-12）を1とします。
\phzu
$f=\qty{2.45e9}{Hz}$、$c=\qty{3.00e8}{m/s}$、波長 $\lambda$。
\phriyu
③の $c=f\lambda$。定在波の腹と腹の間は、半波長（波-6 ②）。
\phtateru
\[ \lambda=\frac{c}{f},\qquad \text{腹と腹の間}=\frac{\lambda}{2} \]
\phtoku
\[ \lambda=\frac{3.00\times10^8}{2.45\times10^9}\approx0.122\ \mathrm{m}=12.2\ \mathrm{cm},\qquad \frac{\lambda}{2}\approx6.1\ \mathrm{cm} \]
\phtashikame
単位：(m/s)÷(1/s)＝m。大きさ：レンジの庫内（数十 cm）に数波長が入る大きさ。皿を回さずに温めると約 \qty{6}{cm} おきにとけむらができる（波-6）。極端な場合：周波数を2倍にすると、波長は半分。
\kotae 波長 約 \qty{12.2}{cm}、腹と腹の間 約 \qty{6.1}{cm}
\end{phreidai}
\begin{jibun}
Wi-Fi の \qty{5.0}{GHz} の電波の波長を求めましょう。
\bkjans{$\lambda=3.00\times10^8\div(5.0\times10^9)=0.060$\,m。\qty{6.0}{cm}}
\end{jibun}

\begin{phreidai}[ラジオ・Wi-Fi の電波の波長]
\label{ph:電-15:ex2}FM ラジオの \qty{80}{MHz}、AM ラジオの \qty{1000}{kHz}、Wi-Fi の \qty{2.4}{GHz} の電波の波長を求めましょう。
\phbamen
②の表の電波です（波-14 ⑤の表と同じ数）。
\phmushi
空気の屈折率を1とします。
\phzu
$f_{\mathrm{FM}}=\qty{8.0e7}{Hz}$、$f_{\mathrm{AM}}=\qty{1.0e6}{Hz}$、$f_{\mathrm{Wi}}=\qty{2.4e9}{Hz}$。
\phriyu
どれも電磁波なので、同じ $c$ で $\lambda=c/f$。
\phtateru
\[ \lambda=\frac{c}{f} \]
\phtoku
\[ \lambda_{\mathrm{FM}}=\frac{3.00\times10^8}{8.0\times10^7}=3.75\ \mathrm{m},\qquad \lambda_{\mathrm{AM}}=\frac{3.00\times10^8}{1.0\times10^6}=300\ \mathrm{m} \]
\[ \lambda_{\mathrm{Wi}}=\frac{3.00\times10^8}{2.4\times10^9}=0.125\ \mathrm{m} \]
\phtashikame
単位：m。大きさ：AM の \qty{300}{m} は、高い鉄塔そのものをアンテナにする大きさ。Wi-Fi の \qty{12.5}{cm} は、手のひらの大きさ（波-14 ⑤）。周波数が2400倍違うと、波長も2400分の1。
\kotae FM 約 \qty{3.75}{m}、AM \qty{300}{m}、Wi-Fi 約 \qty{12.5}{cm}
\end{phreidai}
\begin{jibun}
交通系 IC カードの \qty{13.56}{MHz} と、地上デジタルテレビの \qty{600}{MHz} の電波の波長を求めましょう。
\bkjans{$3.00\times10^8\div(13.56\times10^6)\approx\qty{22}{m}$。$3.00\times10^8\div(6.0\times10^8)=\qty{0.50}{m}$}
\end{jibun}

\begin{phreidai}[太陽の光の電場と磁場]
\label{ph:電-15:ex3}大気の外で、太陽の光に垂直な \qty{1}{m^2} の面が受ける光の強さは \qty{1361}{W/m^2} です。太陽の光を1つの正弦波の電磁波とみなして、電場の振幅と磁場の振幅を求めましょう。
\phbamen
宇宙で太陽の光を受ける人工衛星を考えます。\qty{1361}{W/m^2} は「太陽定数」とよばれる値です。
\phmushi
いろいろな色（波長）がまざった光を、同じ強さの1つの正弦波とみなします（振幅の目安を求める）。
\phzu
$I=\qty{1361}{W/m^2}$、電場の振幅 $E_0$、磁場の振幅 $B_0$。
\phriyu
⑤の $I=\tfrac12c\varepsilon_0E_0^2$ を $E_0$ について解く。磁場は④の $E=cB$。
\phtateru
\[ E_0=\sqrt{\frac{2I}{c\varepsilon_0}},\qquad B_0=\frac{E_0}{c} \]
\phtoku
\[ E_0=\sqrt{\frac{2\times1361}{3.00\times10^8\times8.85\times10^{-12}}}\approx1.0\times10^3\ \mathrm{V/m} \]
\[ B_0=\frac{1.01\times10^3}{3.00\times10^8}\approx3.4\times10^{-6}\ \mathrm{T} \]
\phtashikame
単位：$\sqrt{(\mathrm{W/m^2})\div((\mathrm{m/s})(\mathrm{F/m}))}$＝V/m。大きさ：電場は \qty{1}{m} あたり約 \qty{1000}{V} と、意外に大きい。磁場は地球の磁場の約15分の1。極端な場合：強さが4倍なら、振幅は2倍。
\kotae 電場の振幅 約 $1.0\times10^3$\,V/m、磁場の振幅 約 $3.4\times10^{-6}$\,T
\end{phreidai}
\begin{jibun}
出力 \qty{1.0}{mW} のレーザー光を、\qty{1.0}{mm^2} の面に集めました。強さと、電場の振幅を求めましょう。
\bkjans{$I=1.0\times10^{-3}\div(1.0\times10^{-6})$ $=\qty{1000}{W/m^2}$。$E_0^2=2\times1000$ $\div(3.00\times10^8\times8.85\times10^{-12})$ $\approx7.5\times10^5$ から、$E_0\approx8.7\times10^2$\,V/m}
\end{jibun}

\begin{phreidai}[宇宙ヨットの帆を押す力]
\label{ph:電-15:ex4}差し渡し \qty{20}{m} の正方形の帆（面積 \qty{200}{m^2}）が、太陽の光（\qty{1361}{W/m^2}）を垂直に受けて、全部はね返すとします。帆が受ける放射圧と力を求めましょう。
\phbamen
⑤の「身近な例では」の IKAROS の帆です（大きさは IKAROS と同じ、はね返す割合は仮に全部とします）。
\phmushi
帆の傾き、はね返さずに吸う分、帆のしわを無視します。帆は太陽の光に垂直とします。
\phzu
$I=\qty{1361}{W/m^2}$、面積 $A=\qty{200}{m^2}$、放射圧 $p$、力 $F$。
\phriyu
全部はね返すときの放射圧は⑤の $p=2I/c$。力は圧力×面積。
\phtateru
\[ p=\frac{2I}{c},\qquad F=pA \]
\phtoku
\[ p=\frac{2\times1361}{3.00\times10^8}\approx9.1\times10^{-6}\ \mathrm{Pa},\qquad F=9.07\times10^{-6}\times200\approx1.8\times10^{-3}\ \mathrm{N} \]
\phtashikame
単位：(W/m²)÷(m/s)＝(J/(s·m²))×(s/m)＝N/m²＝Pa。大きさ：\qty{1.8}{mN} は、1円玉（\qty{1}{g}）の重さ（約 \qty{10}{mN}）の5分の1ほど。宇宙では空気の抵抗がないので、この力でも何か月もかけて速さを増やせる。極端な場合：全部吸う黒い帆なら半分。
\kotae 放射圧 約 $9.1\times10^{-6}$\,Pa、力 約 $1.8\times10^{-3}$\,N
\end{phreidai}
\begin{jibun}
太陽の光（\qty{1361}{W/m^2}）を全部吸う、\qty{1.0}{m^2} の黒い板が受ける力を求めましょう。
\bkjans{$p=1361\div(3.00\times10^8)$ $\approx4.5\times10^{-6}$\,Pa。力 約 $4.5\times10^{-6}$\,N}
\end{jibun}

\begin{phreidai}[平面波の磁場の強さ]
\label{ph:電-15:ex5}真空の中を進む平面波の電場の振幅が \qty{1.0}{V/m} です。磁場の強さの振幅と、強さ（平均のポインティングベクトル）を求めましょう。
\phbamen
放送局から離れた所で受ける電波を、平面波とみなします。
\phmushi
アンテナの近くの複雑な場を考えず、平らな波面で進む平面波とします。
\phzu
$E_0=\qty{1.0}{V/m}$、$\eta_0=\qty{377}{\ohm}$、磁場の強さの振幅 $H_0$、強さ $I$。
\phriyu
⑦の $\eta_0=E/H$ と、強さ $I=E_0^2/(2\eta_0)$。
\phtateru
\[ H_0=\frac{E_0}{\eta_0},\qquad I=\frac{E_0^2}{2\eta_0} \]
\phtoku
\[ H_0=\frac{1.0}{377}\approx2.65\times10^{-3}\ \mathrm{A/m},\qquad I=\frac{1.0^2}{2\times377}\approx1.3\times10^{-3}\ \mathrm{W/m^2} \]
\phtashikame
単位：(V/m)÷Ω＝A/m、(V/m)²÷Ω＝W/m²。大きさ：$B_0=\mu_0H_0\approx3.3\times10^{-9}$\,T で、④の $E/c$ と一致。極端な場合：電場が10倍なら、強さは100倍。
\kotae $H_0\approx\qty{2.65}{mA/m}$、$I\approx\qty{1.3}{mW/m^2}$
\end{phreidai}
\begin{jibun}
電場の振幅が \qty{10}{V/m} の平面波の、磁場の強さの振幅と強さを求めましょう。
\bkjans{$H_0=10\div377\approx2.7\times10^{-2}$\,A/m。$I=10^2\div(2\times377)\approx0.13$\,W/m²}
\end{jibun}

\begin{phreidai}[電子と陽子が出す電磁波]
\label{ph:電-15:ex6}電場 \qty{1.0e6}{V/m} の中で、電子と陽子がそれぞれ電場の力だけで加速されています。それぞれが1秒に出す電磁波のエネルギー（仕事率）を求め、比べましょう。電子の質量は $9.11\times10^{-31}$\,kg、陽子の質量は $1.67\times10^{-27}$\,kg とします。
\phbamen
加速器や、雷の強い電場の中の粒を考えます。
\phmushi
重力、ほかの粒とのぶつかり、速さが光に近づく効果（$v\ll c$ とする）を無視します。
\phzu
$E=\qty{1.0e6}{V/m}$、$q=e$、電子の質量 $m_{\mathrm e}$、陽子の質量 $m_{\mathrm p}$、加速度 $a$、仕事率 $P$。
\phriyu
加速度は運動方程式 $ma=eE$（力-2）から。電磁波の仕事率は⑧のラーモアの式。
\phtateru
\[ a=\frac{eE}{m},\qquad P=\frac{e^2a^2}{6\pi\varepsilon_0c^3} \]
\phtoku
電子：$a_{\mathrm e}=1.60\times10^{-19}\times1.0\times10^6\div(9.11\times10^{-31})\approx1.76\times10^{17}$\,m/s²。
\[ P_{\mathrm e}=\frac{(1.60\times10^{-19})^2\times(1.76\times10^{17})^2}{6\pi\times8.85\times10^{-12}\times(3.00\times10^8)^3}\approx1.8\times10^{-19}\ \mathrm{W} \]
陽子：$a_{\mathrm p}\approx9.6\times10^{13}$\,m/s²、$P_{\mathrm p}\approx5.2\times10^{-26}$\,W。比は $(m_{\mathrm p}/m_{\mathrm e})^2\approx3.4\times10^6$。
\phtashikame
単位：C²(m/s²)²÷((F/m)(m/s)³)＝(C²/F)/s＝J/s＝W。大きさ：同じ力でも、軽い電子は加速度が約1800倍で、出す電磁波は約340万倍。光や電波を出すのは、主に電子。極端な場合：$E=0$ なら加速しないので0。
\kotae 電子 約 $1.8\times10^{-19}$\,W、陽子 約 $5.2\times10^{-26}$\,W（電子が約 $3.4\times10^6$ 倍）
\end{phreidai}
\begin{jibun}
電場を3倍の \qty{3.0e6}{V/m} にすると、電子が出す電磁波の仕事率は何倍になりますか。いくらになりますか。
\bkjans{加速度が3倍なので9倍。$1.75\times10^{-19}\times9\approx1.6\times10^{-18}$\,W}
\end{jibun}

\begin{phreidai}[空の青と、ゆれる電子]
\label{ph:電-15:ex7}空気の分子の中の電子がゆすられて出す光の強さは、ゆすられる振幅が色によらないとき、角周波数の4乗に比例します。青い光（\qty{450}{nm}）・緑の光（\qty{550}{nm}）は、赤い光（\qty{700}{nm}）の何倍強く散らされますか。
\phbamen
⑧の空の青さの場面です（波-13 の例題1と同じ数）。
\phmushi
太陽の光の色ごとの強さの違い、目の感じ方、ちりや水の粒の散乱を無視します。電子の固有振動数は可視光よりずっと高いとします（振幅が色によらない）。
\phzu
$\lambda_{\text{青}}=\qty{450}{nm}$、$\lambda_{\text{緑}}=\qty{550}{nm}$、$\lambda_{\text{赤}}=\qty{700}{nm}$。
\phriyu
ラーモアの式で仕事率は加速度の2乗、加速度の振幅は $\omega^2x_0$ なので $\omega^4x_0^2$。$\omega=2\pi c/\lambda$ なので $1/\lambda^4$ に比例。
\phtateru
\[ \frac{P_{\text{青}}}{P_{\text{赤}}}=\Bigl(\frac{\lambda_{\text{赤}}}{\lambda_{\text{青}}}\Bigr)^4 \]
\phtoku
\[ \Bigl(\frac{700}{450}\Bigr)^4\approx5.9,\qquad \Bigl(\frac{700}{550}\Bigr)^4\approx2.6 \]
\phtashikame
単位なし。大きさ：波-13 の例題1と同じ答え。波-13 では「加速度の2乗」を結果として使ったが、ここでは⑧の電場の折れ曲がりから導いた式で計算した。極端な場合：同じ波長なら1。
\kotae 青は約5.9倍、緑は約2.6倍
\end{phreidai}
\begin{jibun}
紫の光（\qty{400}{nm}）は、赤い光（\qty{650}{nm}）の何倍強く散らされますか。
\bkjans{$(650\div400)^4=1.625^4\approx7.0$。約7.0倍}
\end{jibun}

\begin{phreidai}[ガラスと水の表面の反射率]
\label{ph:電-15:ex8}空気から、ガラス（屈折率1.5）と水（屈折率1.333）の表面に、光が垂直に当たります。電場の反射係数と反射率を求めましょう。はね返った光の山と谷は入れかわりますか。
\phbamen
窓ガラスや水面の照り返しです（波-7 ⑨・波-13 の図と同じ数）。
\phmushi
ガラスと水は光を吸わないとし、空気の屈折率を1とします。裏の面の反射は考えません。
\phzu
$n_1=1$、ガラス $n_2=1.5$、水 $n_2=1.333$。反射係数 $r_E$、反射率 $R$。
\phriyu
⑨の、境目の電場と磁場のつながり方から導いた $r_E=(n_1-n_2)/(n_1+n_2)$、$R=r_E^{\,2}$。
\phtateru
\[ r_E=\frac{n_1-n_2}{n_1+n_2},\qquad R=r_E^{\,2} \]
\phtoku
ガラス：$r_E=(1-1.5)/(1+1.5)=-0.20$、$R=0.040$。水：$r_E=(1-1.333)/(1+1.333)\approx-0.143$、$R\approx0.020$。
\phtashikame
単位なし。大きさ：どちらも数 \% で、波-7 ⑨・波-13 の図の真上からの値と同じ。$r_E$ が負なので、どちらも山と谷が入れかわる（⑨）。極端な場合：$n_2=1$（境目がない）なら0。
\kotae ガラス $r_E=-0.20$、4.0\,\%。水 $r_E\approx-0.14$、約2.0\,\%。どちらも山と谷が入れかわる
\end{phreidai}
\begin{jibun}
空気からダイヤモンド（屈折率2.42）に、光が垂直に当たるときの反射率を求めましょう。
\bkjans{$r_E=(1-2.42)\div(1+2.42)\approx-0.415$。$R\approx0.17$。約17\,\%}
\end{jibun}

\begin{phreidai}[ガラスと水のブルースター角]
\label{ph:電-15:ex9}空気からガラス（屈折率1.5）と水（屈折率1.333）へ光が入るときの、ブルースター角を求めましょう。その入射角で、反射光と屈折光のなす角が $90°$ になることを確かめましょう。
\phbamen
波-13 ⑥で、偏光サングラスやカメラのフィルターで照り返しを消した場面です。
\phmushi
空気の屈折率を1とし、表面は平らとします。
\phzu
$n_1=1$、$n_2$、ブルースター角 $\theta_{\mathrm B}$、屈折角 $\theta_2$。
\phriyu
⑨で、$r_\parallel=0$ になる条件から $\tan\theta_{\mathrm B}=n_2/n_1$ と $\theta_1+\theta_2=90°$ が出た。
\phtateru
\[ \tan\theta_{\mathrm B}=\frac{n_2}{n_1},\qquad \sin\theta_2=\frac{n_1\sin\theta_{\mathrm B}}{n_2} \]
\phtoku
ガラス：$\tan\theta_{\mathrm B}=1.5$ から $\theta_{\mathrm B}\approx56.3°$。$\sin\theta_2=\sin56.3°\div1.5\approx0.555$、$\theta_2\approx33.7°$。和は $90.0°$。\par
水：$\tan\theta_{\mathrm B}=1.333$ から $\theta_{\mathrm B}\approx53.1°$、$\theta_2\approx36.9°$。和は $90.0°$。
\phtashikame
単位：度。大きさ：波-13 の図の $56.3°$ と同じ。反射光と屈折光が直角なので、屈折光の中で紙面の中にゆれる電子のゆれの向きが、反射光の向きと一致し、⑧のとおりその向きには光を出さない。極端な場合：$n_2\to1$ なら $45°$。
\kotae ガラス 約 $56.3°$、水 約 $53.1°$（どちらも屈折角との和が $90°$）
\end{phreidai}
\begin{jibun}
空気からダイヤモンド（屈折率2.42）へ光が入るときのブルースター角を求めましょう。
\bkjans{$\tan\theta_{\mathrm B}=2.42$ から約 $67.5°$}
\end{jibun}

\begin{phreidai}[水の比誘電率と屈折率]
\label{ph:電-15:ex10}水の比誘電率は、コンデンサーではかると約80（\qty{20}{\celsius}）です。⑥の $n=\sqrt{\varepsilon_{\mathrm r}}$ で屈折率を計算し、光の屈折率1.333と比べましょう。光に対する水の比誘電率はいくらですか。
\phbamen
⑨の、水の比誘電率と屈折率の場面です。
\phmushi
水は磁石につかない（$\mu_{\mathrm r}=1$）とします。光は吸われないとします。
\phzu
ゆっくり変わる電場での比誘電率 $\varepsilon_{\mathrm r}=80$、光の屈折率 $n=1.333$。
\phriyu
⑥の $n=\sqrt{\varepsilon_{\mathrm r}\mu_{\mathrm r}}$。逆に、$\varepsilon_{\mathrm r}=n^2$ で光に対する比誘電率がわかる。
\phtateru
\[ n=\sqrt{\varepsilon_{\mathrm r}},\qquad \varepsilon_{\mathrm r}=n^2 \]
\phtoku
$\sqrt{80}\approx8.9$。光の屈折率1.333とは大きく違う。光に対しては $\varepsilon_{\mathrm r}=1.333^2\approx1.78$。
\phtashikame
単位なし。大きさ：光の周波数（$10^{14}$\,Hz 以上）では、水の分子の回転が電場の速いゆれについていけず、分極が小さくなる（⑨）。マイクロ波（\qty{2.45}{GHz}）では分子の回転が少しおくれてついていくので、そのおくれの分でエネルギーを吸い、水が温まる（電子レンジ）。極端な場合：ゆっくり変わる電場ほど、比誘電率は80に近い。
\kotae $\sqrt{80}\approx8.9$（光の屈折率1.333と合わない）。光に対する比誘電率は約1.78
\end{phreidai}
\begin{jibun}
ガラス（光の屈折率1.5）の、光に対する比誘電率を求めましょう。
\bkjans{$1.5^2=2.25$}
\end{jibun}

\begin{phreidai}[銅の表皮深さ]
\label{ph:電-15:ex11}銅（導電率 \qty{5.8e7}{S/m}、$\mu_{\mathrm r}=1$）の表皮深さを、\qty{50}{Hz}・\qty{1}{MHz}・\qty{2.45}{GHz} で求めましょう。
\phbamen
家の電気（\qty{50}{Hz}）、ラジオ（\qty{1}{MHz}）、電子レンジ（\qty{2.45}{GHz}）の電流が、銅の表面のどれくらいの深さを流れるかを考えます。
\phmushi
銅は一様で、表面は平らとします（表皮深さが銅の厚さよりずっと小さいとき）。
\phzu
$\sigma=\qty{5.8e7}{S/m}$、$\mu=\mu_0$、$\omega=2\pi f$、表皮深さ $\delta$。
\phriyu
⑩（電-8 ⑫）の $\delta=\sqrt{2/(\omega\mu\sigma)}$。
\phtateru
\[ \delta=\sqrt{\frac{2}{2\pi f\mu_0\sigma}}=\frac{1}{\sqrt{\pi f\mu_0\sigma}} \]
\phtoku
$\pi\mu_0\sigma=\pi\times4\pi\times10^{-7}\times5.8\times10^7\approx229$ なので、$\delta=1/\sqrt{229f}$。
\[ 50\ \mathrm{Hz}:\ \delta=\frac{1}{\sqrt{229\times50}}\approx9.3\times10^{-3}\ \mathrm{m},\qquad 1\ \mathrm{MHz}:\ \delta\approx6.6\times10^{-5}\ \mathrm{m} \]
\[ 2.45\ \mathrm{GHz}:\ \delta\approx1.3\times10^{-6}\ \mathrm{m} \]
\phtashikame
単位：$\sqrt{1/((1/\mathrm s)(\mathrm{N/A^2})(\mathrm{S/m}))}$＝m。大きさ：\qty{50}{Hz} の約 \qty{9}{mm} は太い送電線の太さくらいで、ふつうの電線は全体に電流が流れる。\qty{2.45}{GHz} では \qty{1.3}{\micro m} と、髪の毛の太さの50分の1ほどの表面だけ。極端な場合：周波数を4倍にすると、表皮深さは半分。
\kotae \qty{50}{Hz} 約 \qty{9.3}{mm}、\qty{1}{MHz} 約 \qty{66}{\micro m}、\qty{2.45}{GHz} 約 \qty{1.3}{\micro m}
\end{phreidai}
\begin{jibun}
銅の表皮深さを、\qty{60}{Hz} と \qty{13.56}{MHz} で求めましょう。
\bkjans{$1/\sqrt{229\times60}$ $\approx8.5\times10^{-3}$\,m（約 \qty{8.5}{mm}）。$229\times1.356\times10^7\approx3.1\times10^9$ なので $\delta\approx1.8\times10^{-5}$\,m（約 \qty{18}{\micro m}）}
\end{jibun}

\begin{phreidai}[電子レンジの網の穴]
\label{ph:電-15:ex12}電子レンジの窓の網の穴を、直径 \qty{1.0}{mm}、深さ（網の厚さ）\qty{1.0}{mm} の円い穴とします。(1) 穴を通れるいちばん長い波長（遮断波長）を求めましょう。(2) \qty{2.45}{GHz} のマイクロ波の電場は、穴を通りぬけるまでに何分の1に弱まりますか。強さ（エネルギー）では何分の1ですか。
\phbamen
はてな？の電子レンジの窓です。穴の大きさは、仮の値とします。
\phmushi
穴が1つだけのときを考え、周りの穴どうしの影響、網のふちのすきまを無視します。網は電気をよく通す金属とします。
\phzu
穴の直径 $d=\qty{1.0e-3}{m}$、厚さ $L=\qty{1.0e-3}{m}$、遮断波長 $\lambda_{\mathrm c}$。マイクロ波の波長は \qty{12.2}{cm}（例題1）。
\phriyu
⑩の $\lambda_{\mathrm c}\approx1.71d$。マイクロ波の波長は $\lambda_{\mathrm c}$ よりずっと長いので、穴の中で、長さ $\lambda_{\mathrm c}/(2\pi)$ ごとに $1/e$ に弱まる。強さは電場の2乗に比例する（⑤）。
\phtateru
\[ \lambda_{\mathrm c}=1.71\,d,\qquad \frac{E_{\text{出}}}{E_{\text{入}}}=e^{-2\pi L/\lambda_{\mathrm c}},\qquad \frac{I_{\text{出}}}{I_{\text{入}}}=\Bigl(\frac{E_{\text{出}}}{E_{\text{入}}}\Bigr)^2 \]
\phtoku
(1) $\lambda_{\mathrm c}=1.71\times1.0\times10^{-3}\approx1.7\times10^{-3}$\,m（約 \qty{1.7}{mm}）。\par
(2) $2\pi L/\lambda_{\mathrm c}=2\pi\times1.0\div1.71\approx3.68$。$e^{-3.68}\approx0.025$（約40分の1）。強さは $0.025^2\approx6.3\times10^{-4}$（約1600分の1）。
\phtashikame
単位：mm÷mm で、指数の中は単位なし。大きさ：デシベルでは $10\log_{10}1600\approx\qty{32}{dB}$（3-12⑦）。可視光（約 \qty{0.5}{\micro m}）は遮断波長よりずっと短いので、弱まらずに通る。極端な場合：網を2倍の厚さにすると、強さはさらに1600分の1。穴を大きくすると、遮断波長が長くなり、弱め方が小さくなる。
\kotae (1) 約 \qty{1.7}{mm}　(2) 電場は約40分の1、強さは約1600分の1（約 \qty{32}{dB}）
\end{phreidai}
\begin{jibun}
例題12の網の穴の直径を \qty{2.0}{mm} にすると、\qty{2.45}{GHz} の電磁波の強さは何分の1になりますか（厚さは \qty{1.0}{mm}）。
\bkjans{$\lambda_{\mathrm c}\approx3.4$\,mm。$e^{-2\pi\times1.0/3.42}\approx0.16$。強さは $0.16^2\approx0.025$ で約40分の1}
\end{jibun}

\begin{phreidai}[アンテナの長さと、出す電力]
\label{ph:電-15:ex13}(1) \qty{2.4}{GHz} の Wi-Fi の電波を出す、金属の板の上に立てる $\tfrac14$ 波長のアンテナの長さを求めましょう。(2) 放射抵抗 \qty{73}{\ohm} の半波長のアンテナの根もとに、実効値 \qty{1.0}{A} の電流を流すと、何 W の電波が出ますか。
\phbamen
⑪のアンテナです。(1) は波-14 ⑤の例題1と同じ数です。
\phmushi
アンテナの太さ、周りの物、アンテナの中で熱になる分を無視します。
\phzu
$f=\qty{2.4e9}{Hz}$、アンテナの長さ $\lambda/4$。放射抵抗 $R_{\mathrm r}=\qty{73}{\ohm}$、電流の実効値 \qty{1.0}{A}。
\phriyu
(1) $\lambda=c/f$。(2) ⑪の $P=R_{\mathrm r}\times$（電流の実効値）$^2$。
\phtateru
\[ \frac{\lambda}{4}=\frac{c}{4f},\qquad P=R_{\mathrm r}\times(1.0)^2 \]
\phtoku
(1) $\lambda=3.00\times10^8\div(2.4\times10^9)=0.125$\,m、$\lambda/4\approx0.031$\,m（約 \qty{3.1}{cm}）。(2) $P=73\times1.0^2=\qty{73}{W}$。
\phtashikame
単位：m、Ω×A²＝W。大きさ：(1) はルーターのアンテナの長さと合う。(2) は放送局なら小さく、スマホ（1 W 以下）よりずっと大きい。極端な場合：電流が半分なら電力は4分の1。
\kotae (1) 約 \qty{3.1}{cm}　(2) \qty{73}{W}
\end{phreidai}
\begin{jibun}
FM ラジオ（\qty{80}{MHz}）の半波長のアンテナの長さを求めましょう。このアンテナに実効値 \qty{0.50}{A} を流すと、何 W の電波が出ますか（放射抵抗 \qty{73}{\ohm}）。
\bkjans{$\lambda/2=3.75\div2\approx\qty{1.9}{m}$。$P=73\times0.50^2\approx\qty{18}{W}$}
\end{jibun}

\begin{phreidai}[電波の強さと、電波防護指針]
\label{ph:電-15:ex14}\qty{1.0}{W} の電波を、どの向きにも同じ強さで出すアンテナがあるとします。\qty{10}{m} 離れた所の電波の強さを求め、電波防護指針の値（\qty{1.5}{GHz} 以上で \qty{10}{W/m^2}）と比べましょう。
\phbamen
⑫の電波の安全の場面です。
\phmushi
地面や建物での反射、空気での吸収を無視します。アンテナを点とみなします（どの向きにも同じ強さの、仮のアンテナ）。
\phzu
仕事率 $P=\qty{1.0}{W}$、距離 $r=\qty{10}{m}$、強さ $I$。
\phriyu
とちゅうで吸われなければ、どの球面を通るエネルギーも同じ（逆2乗の法則：波-7 ②）。
\phtateru
\[ I=\frac{P}{4\pi r^2} \]
\phtoku
$I=1.0\div(4\pi\times10^2)\approx8.0\times10^{-4}$\,W/m²。指針の \qty{10}{W/m^2} の約 $1.3\times10^4$ 分の1。
\phtashikame
単位：W÷m²。大きさ：指針の値を超えるのは、$1.0\div(4\pi r^2)=10$ から $r\approx\qty{9}{cm}$ より近い所だけ。距離が2倍で強さは4分の1。
\kotae 約 $8.0\times10^{-4}$\,W/m²（指針の値の約1万3千分の1）
\end{phreidai}
\begin{jibun}
\qty{0.20}{W} の電波をどの向きにも同じ強さで出すアンテナから \qty{1.0}{m} 離れた所の強さを求め、\qty{10}{W/m^2} と比べましょう。
\bkjans{$0.20\div(4\pi\times1.0^2)\approx0.016$\,W/m²。約630分の1}
\end{jibun}

%% ---------------- まちがえやすいところ ----------------
\bksection{tsumazuki}{まちがえやすいところをチェック}
\begin{tsumazuki}{周波数が高い電磁波ほど速いと考える}
\bkng{X 線は周波数が高いので、電波より速く進む。}
\bkok{真空の中では、どの電磁波も同じ速さ $c$。周波数が高いと、波長が短くなるだけ。}
$c=f\lambda$ で、$c$ は決まった値です。周波数が高くなれば、そのぶん波長が短くなります。物の中では、周波数で速さが少し変わります（分散：⑨）。
\end{tsumazuki}
\begin{tsumazuki}{電磁波の磁場は小さいので、エネルギーも小さいと考える}
\bkng{$B=E/c$ はとても小さいので、電磁波のエネルギーはほとんど電場が運ぶ。}
\bkok{エネルギー密度で比べると、磁場の $B^2/(2\mu_0)$ と電場の $\tfrac12\varepsilon_0E^2$ はちょうど等しい。}
磁場の数（T）は小さく見えても、エネルギー密度の式の $1/\mu_0$ が大きいので、電場と同じだけのエネルギーを運びます（⑤）。単位の違う量を、数の大きさだけで比べてはいけません。
\end{tsumazuki}
\begin{tsumazuki}{金属の網は、穴より小さい波を止めると考える}
\bkng{網の穴より波長が小さいマイクロ波は通りぬけ、光は止まる。}
\bkok{波長が遮断波長（穴の直径の約1.7倍）より長い電磁波が止められる。光は波長が短いので通る。}
穴を通れるかどうかは、「波長が穴に入るか」で決まります。マイクロ波の波長（\qty{12.2}{cm}）は穴よりずっと大きく、穴の中で定在波を作れません。
\end{tsumazuki}
\begin{tsumazuki}{反射で山と谷が入れかわる向きを逆に覚える}
\bkng{屈折率の小さい物との境（ガラス→空気）での反射で、山と谷が入れかわる。}
\bkok{入れかわるのは、屈折率の大きい物との境（空気→ガラス）での反射。$r_E=(n_1-n_2)/(n_1+n_2)$ が負になるとき。}
迷ったら、⑨の式に数を入れて、$r_E$ の符号を見ます。$n_1=1$、$n_2=1.5$ なら $-0.2$（入れかわる）、逆なら $+0.2$（そのまま）です。
\end{tsumazuki}

%% ---------------- 工学ではこう使う ----------------
\bksection{phgear}{工学ではこう使う}
\begin{phkougaku}[電磁波で、送る・温める・はかる]
\phboxhead{① 通信と放送}
スマホ・テレビ・GPS・ETC など、たくさんの無線が同時に使われています。電波は限られた資源なので、周波数ごとに使いみちと使う人が法律で決められ、出してよい強さや、ほかの周波数へのもれも決められています。

\phboxhead{② 電磁波で温める}
電子レンジは、水の分子がマイクロ波の電場で回ろうとして、少しおくれてついていくときに失うエネルギーで、食べ物を温めます（例題10）。工場では、同じしくみで、木材やゴムを中から乾かしたり、固めたりします。

\phboxhead{③ 電磁波ではかる}
レーダーは、電磁波を出して、はね返ってくるまでの時間と周波数のずれ（ドップラー効果：波-8 ⑦）から、距離と速さをはかります（波-14 ⑦）。天気予報の雨雲レーダー、車の衝突防止、空港の管制に使われます。電磁波の反射を使って、地中や壁の中を調べる装置もあります。

\begin{phketa}[電磁波の数]
\phketaline{真空の中の光の速さ}{\qty{299792458}{m/s}（決めた値）}
\phketaline{真空の波のインピーダンス}{約 \qty{377}{\ohm}}
\phketaline{電子レンジのマイクロ波}{\qty{2.45}{GHz}・\qty{12.2}{cm}}
\phketaline{太陽の光の強さ（大気の外）}{\qty{1361}{W/m^2}}
\phketaline{半波長のアンテナの放射抵抗}{約 \qty{73}{\ohm}}
\phboxhead{答えが桁で合っているか}
アンテナの大きさは波長の半分か4分の1、金属の穴は波長よりずっと小さいと電磁波を通さない。波長を $c/f$ で出してから、物の大きさと比べると、答えの桁が確かめられます。
\end{phketa}
\end{phkougaku}
\phstation{電波工学,光工学,通信工学,EMC}{電磁波、平面波、波動インピーダンス、ポインティングベクトル、放射圧、ダイポールアンテナ、放射抵抗、放射パターン、フレネルの式、表皮効果、遮へい、導波管、遮断波長}

\begin{column}{ヘルツが確かめた、目に見えない光}
マクスウェルが電磁波を予言してから20年あまりたった1888年ごろ、ドイツのヘルツは、2つの金属の玉のすきまに火花を飛ばして、電気のゆれを作りました。すると、離れた所に置いた輪の形の電線のすきまにも、小さな火花が飛びました。ヘルツは、この「目に見えない波」が、金属の板ではね返ること、定在波を作ること、曲がって進むことを確かめ、それが光と同じ性質をもつ電磁波であることを示しました。周波数の単位ヘルツ（Hz）は、この人の名前です。その後、マルコーニたちが電磁波を通信に使い始め（1909年に、ブラウンとともにノーベル物理学賞を受けた）、ラジオ・テレビ・スマホの時代へとつながりました。
\end{column}

\begin{phsaki}{電波工学・光工学と、光の粒}
この章では、電磁波の基本の性質と、簡単な形のアンテナ・穴・境目だけを調べました。電波工学の専門の本では、マクスウェル方程式を、電位とベクトルポテンシャル（電-6）の、おくれて伝わる形（遅延ポテンシャル）に書きかえて、いろいろな形のアンテナが出す電波や、金属の管（導波管）の中を進む電磁波を計算します。光工学の本では、光ファイバの中の光、レーザーのしくみ、薄い膜の重ね方（波-13 ④）を学びます。

一方、光を当てると金属から電子が飛び出す光電効果は、電磁波の式だけでは説明できず、光がエネルギーのつぶ（光子）でもあることを示しています。物理シリーズの『原子』（原-2）では、⑤の「運動量はエネルギー÷$c$」を、光子1つ1つに当てはめます。
\end{phsaki}

%% ---------------- 章末問題 ----------------
\bkendproblems
\bkproblemgroup{基本}{3}
\bkproblem AM ラジオの \qty{594}{kHz} と、FM ラジオの \qty{90}{MHz} の電波の波長を求めましょう（\nobreak\bkreidairef{2}\nobreak）。
\bkpans{$3.00\times10^8\div(5.94\times10^5)\approx\qty{505}{m}$（約 \qty{510}{m}）。$3.00\times10^8\div(9.0\times10^7)\approx\qty{3.3}{m}$}
\bkproblem 波長 \qty{6.0}{cm} の電磁波の周波数を求めましょう（\nobreak\bkreidairef{1}\nobreak）。
\bkpans{$f=3.00\times10^8\div0.060=5.0\times10^9$\,Hz（\qty{5.0}{GHz}）}
\bkproblem 電場の振幅が \qty{300}{V/m} の電磁波の、磁場の振幅を求めましょう（④・\nobreak\bkreidairef{3}\nobreak）。
\bkpans{$B_0=300\div(3.00\times10^8)=1.0\times10^{-6}$\,T}
\bkproblem 電場の振幅が \qty{100}{V/m} の電磁波の強さを求めましょう（\nobreak\bkreidairef{3}\nobreak・\nobreak\bkreidairef{5}\nobreak）。
\bkpans{$I=\tfrac12\times3.00\times10^8$ $\times8.85\times10^{-12}\times100^2$ $\approx\qty{13}{W/m^2}$}
\bkproblem 空気から屈折率1.6のプラスチックのレンズに、光が垂直に当たるときの反射率を求めましょう（\nobreak\bkreidairef{8}\nobreak）。
\bkpans{$r_E=(1-1.6)\div(1+1.6)\approx-0.23$。$R\approx0.053$。約5.3\,\%}
\bkproblem 青い光（\qty{470}{nm}）は、赤い光（\qty{680}{nm}）の何倍強く、空気の分子に散らされますか（\nobreak\bkreidairef{7}\nobreak）。
\bkpans{$(680\div470)^4\approx4.4$。約4.4倍}
\bkproblemgroup{チャレンジ}{5}
\bkproblem 電場 \qty{2.0e5}{V/m} の中で加速される電子が出す電磁波の仕事率を求めましょう。この電子が1秒に出すエネルギーを、電子ボルト（\qty{1}{eV}$=1.60\times10^{-19}$\,J）で表しましょう（\nobreak\bkreidairef{6}\nobreak）。
\bkpans{$a=1.60\times10^{-19}\times2.0\times10^5$ $\div(9.11\times10^{-31})$ $\approx3.5\times10^{16}$\,m/s²。$P=e^2a^2/(6\pi\varepsilon_0c^3)$ $\approx7.0\times10^{-21}$\,W。1秒に約0.044\,eV}
\bkproblem 空気から屈折率1.6のガラスに入る光のブルースター角と、そのときの屈折角を求めましょう。この角で反射した光は、どんな光になっていますか（\nobreak\bkreidairef{9}\nobreak）。
\bkpans{$\tan\theta_{\mathrm B}=1.6$ から約 $58.0°$、屈折角は約 $32.0°$（和が $90°$）。反射光は紙面に垂直にゆれる光だけ（偏光）になっている}
\bkproblem アルミニウム（導電率 \qty{3.8e7}{S/m} とする）の \qty{2.45}{GHz} での表皮深さを求めましょう。台所用のアルミはくの厚さを \qty{11}{\micro m} とすると、表皮深さの何倍ですか（\nobreak\bkreidairef{11}\nobreak）。
\bkpans{$\pi f\mu_0\sigma=\pi\times2.45\times10^9$ $\times4\pi\times10^{-7}\times3.8\times10^7$ $\approx3.7\times10^{11}$ なので、$\delta\approx1.6\times10^{-6}$\,m（約 \qty{1.6}{\micro m}）。$11\div1.65\approx6.7$ 倍。電磁波は、はくの中をほとんど通れない}
\bkproblem 電子レンジの網の穴の直径を \qty{1.0}{mm} とします。網の厚さを何 mm にすれば、\qty{2.45}{GHz} の電磁波の強さを1万分の1（\qty{40}{dB}）以下にできますか（\nobreak\bkreidairef{12}\nobreak・5-14）。
\bkpans{強さは $e^{-2\times2\pi L/\lambda_{\mathrm c}}$。$e^{-4\pi L/1.71}\leqq10^{-4}$ から $4\pi L/1.71\geqq\ln10^4\approx9.21$、$L\geqq1.25$\,mm。約 \qty{1.3}{mm} 以上}

\begin{furikaeri}
\begin{checklist}
\item 電波・光・X 線などが、波長（周波数）の違う同じ電磁波であることを知り、$c=f\lambda$ で波長を計算できる。\bkfurref{ph:電-15:sec1}
\item 電磁波が電場と磁場の直角にゆれる横波で、$E=cB$ の関係をもち、エネルギーと運動量を運ぶことを計算で確かめられる。\bkfurref{ph:電-15:sec4}
\item ゆれる電荷が出す電磁波（ラーモアの式）、物の表面での反射率、金属の網が電波を通さないわけを、マクスウェル方程式から説明できる。\bkfurref{ph:電-15:sec8}
\end{checklist}
\end{furikaeri}
\bkfukidashi{mina}{電気と磁気の式から、光の速さも、空の青さも、電子レンジの網のわけも出てきたね。最後の電-16 では、電気を作って家まで届けるしくみを、まとめて読んでいくよ。}
